\documentclass[11pt]{article}

\usepackage[margin=1in]{geometry}
\usepackage{amsmath,amssymb,amsthm}
\usepackage{graphicx}
\usepackage{booktabs}
\usepackage{hyperref}
% [numbers] switches natbib from its default author-year mode to numeric
% citations; unsrtnat then numbers the bibliography in order of first citation.
% \citet, \citep and \citeauthor all keep working.
\usepackage[numbers]{natbib}
\usepackage[table]{xcolor}

% Sections and subsections only; there are no subsubsections.
\setcounter{tocdepth}{2}

% ---------------------------------------------------------------------
% Provenance tags for parameter values (Section~\ref{sec:params}).
% LIT   -- derived from published CWD field or modelling data
% DEF   -- a unit-defining normalization; exact by construction
% CAL   -- back-solved from a calibration target (a target R0, a term balance)
% GUESS -- expert judgement with no supporting citation for this study region
% Rows carrying a GUESS value are shaded so they can be found at a glance.
% ---------------------------------------------------------------------
\definecolor{guessbg}{rgb}{1.00,0.94,0.80}
\newcommand{\tagLIT}{{\footnotesize\textsf{\textcolor{green!40!black}{LIT}}}}
\newcommand{\tagDEF}{{\footnotesize\textsf{\textcolor{blue!60!black}{DEF}}}}
\newcommand{\tagCAL}{{\footnotesize\textsf{\textcolor{orange!75!black}{CAL}}}}
\newcommand{\tagGUESS}{{\footnotesize\textsf{\textbf{\textcolor{red!70!black}{GUESS}}}}}
\newcommand{\tagNUM}{{\footnotesize\textsf{\textcolor{purple!70!black}{NUMERICAL}}}}
\newcommand{\guessrow}{\rowcolor{guessbg}}

\newcommand{\Grad}{\nabla_\Gamma}
\newcommand{\Div}{\nabla_\Gamma\cdot}
\newcommand{\Vad}{\mathcal{V}_{\mathrm{ad}}}
\newcommand{\Ham}{\mathcal{H}}
\newcommand{\Lag}{\mathcal{L}}
\title{A Reaction--Diffusion Model for Chronic Wasting Disease Spread in Deer on a Topographically Realistic Landscape, with Optimal Culling Control}
\author{}
\date{\today}

\begin{document}
\maketitle

\begin{abstract}
We present a spatially explicit model of chronic wasting disease (CWD) transmission in white-tailed deer, posed on a two-dimensional surface $\Gamma \subset \mathbb{R}^3$ that is fit to real topographic data. The model extends the classical directly-transmitted wildlife disease framework by including an environmental compartment representing prion contamination of the landscape, which allows for indirect transmission in addition to direct deer-to-deer contact. Diffusion of the deer sub-populations is modeled anisotropically, with local diffusivity depending on land cover and terrain slope, and the local deer carrying capacity is informed by land cover classification data. We describe the model, its parameterization from geographic data sources, and a surface finite element method (FEM) suitable for its numerical solution. We then pose a management problem on top of the model: choosing a spatially and temporally varying culling effort, applied to the prion-shedding class, so as to minimize a weighted combination of the shedding burden, the environmental prion load, and the cost of the intervention. We derive the adjoint system characterizing the optimal control, give the discrete adjoint of the implicit--explicit time-stepping scheme, and describe a forward--backward sweep driven by a projected quasi-Newton method.
\end{abstract}

\tableofcontents
\newpage

\section{Introduction}

Chronic wasting disease is a fatal, prion-mediated neurodegenerative disease affecting deer, elk, and other cervids. Unlike many classical wildlife epizootics, CWD is transmitted not only through direct contact between infectious and susceptible animals but also indirectly, through environmental contamination: prions shed by infectious deer persist in soil and vegetation for extended periods and remain capable of infecting susceptible deer long after the shedding animal has died or moved on \citep{Almberg2011}. Capturing this environmental reservoir is essential to realistically modeling CWD spread and to evaluating management strategies such as targeted culling or carcass removal.

Real deer habitats are neither flat nor homogeneous: terrain elevation, slope, and land cover all influence both the local population density a landscape can support and the rate at which deer -- and hence the disease and environmental contamination -- move through it. In this work we pose a four-compartment susceptible--infected--infectious--environmental (S-I-R-W) reaction--diffusion model on a two-dimensional surface $\Gamma$ embedded in $\mathbb{R}^3$, constructed by interpolating digital elevation data. We couple this surface to land-cover classification data so that both the deer carrying capacity $K$ and the diffusion tensor $D$ vary in space according to real geographic features, and describe a surface finite element discretization appropriate for integrating the resulting system on an unstructured triangulated mesh of $\Gamma$.

Sections~\ref{sec:control}--\ref{sec:control-implementation} take the further step of treating the model as a constraint rather than as an end in itself. Given the model, one may ask how a limited culling effort should be distributed over space and time in order to do the most good. We formulate this as an optimal control problem, derive the adjoint equations that characterize its solution, and describe the numerical method used to solve it.

\section{Domain: The Embedded Surface $\Gamma$}
\label{sec:domain}

The spatial domain $\Gamma \subset \mathbb{R}^3$ is a two-dimensional surface constructed by interpolating topographic elevation data obtained from the USGS National Map downloader (\url{https://apps.nationalmap.gov/downloader/}). The resulting digital elevation model (DEM) is triangulated to produce a mesh $\Gamma_h$ (Section~\ref{sec:fem}) whose vertices carry both a spatial position $\mathbf{x} \in \mathbb{R}^3$ and an outward unit surface normal $\mathbf{n}(\mathbf{x})$.

Land cover classification data, obtained from the Multi-Resolution Land Characteristics Consortium viewer (\url{https://www.mrlc.gov/viewer/}), is glued to the elevation surface by aligning shared longitude and latitude coordinates. Each point of $\Gamma$ is thus tagged with a land cover class, which is used below to define both the deer carrying capacity $K$ and the diffusion tensor $D$.

All spatial differential operators appearing in the model are surface (Laplace--Beltrami type) operators: for a scalar field $u$ on $\Gamma$, $\Grad u$ denotes the tangential gradient and $\Div(\mathbf{F})$ the surface divergence of a tangential vector field $\mathbf{F}$.

\section{The Reaction--Diffusion Model}
\label{sec:model}

\subsection{State variables and governing equations}

The deer population is partitioned into three compartments, each a spatially and temporally varying density (deer/km$^2$):
\begin{itemize}
  \item $S(\mathbf{x},t)$ --- \emph{susceptible} deer;
  \item $I(\mathbf{x},t)$ --- \emph{infected} deer, in the incubating, sub-clinical, non-shedding stage;
  \item $R(\mathbf{x},t)$ --- \emph{shedding} deer, actively depositing prions. This class lumps the pre-clinical shedding stage together with the terminal clinical stage; the consequences are taken up in Section~\ref{sec:params}.
\end{itemize}
We write $N = S + I + R$ for the total local deer density. To these we adjoin a fourth field,
\begin{itemize}
  \item $W(\mathbf{x},t)$ --- the \emph{environmental prion load}, the local density of infectious prion material deposited on the landscape through saliva, urine, feces, and decomposing carcasses,
\end{itemize}
which tracks the contaminated-environment reservoir that distinguishes CWD from directly transmitted wildlife diseases; our notation for this compartment follows the environmental-transmission literature on CWD \citep{Almberg2011}.

The governing system on $\Gamma$ is
\begin{align}
  \partial_t S &= \Div\big(D_S\,\Grad S\big) + rS\Big(1-\frac{N}{K}\Big) - \beta_1 S R - \beta_2 S W, \label{eq:S} \\
  \partial_t I &= \Div\big(D_I\,\Grad I\big) + \beta_1 S R + \beta_2 S W - \sigma I, \label{eq:I} \\
  \partial_t R &= \Div\big(D_R\,\Grad R\big) + \sigma I - \alpha R, \label{eq:R} \\
  \partial_t W &= \rho R - \delta W, \label{eq:W}
\end{align}
where the diffusion tensor $D$ is specified in Section~\ref{sec:diffusion} and the carrying capacity $K = K(\mathbf{x})$ in Section~\ref{sec:carrying}. The system is closed with no-flux (homogeneous Neumann) conditions
\begin{equation}
  \big(D_X\,\Grad X\big)\cdot\boldsymbol{\nu} = 0
  \qquad \text{on } \partial\Gamma, \qquad X\in\{S,I,R\},
  \label{eq:bc}
\end{equation}
$\boldsymbol{\nu}$ being the outward conormal, so that deer neither enter nor leave the domain across its boundary, and with prescribed nonnegative initial data.

\subsection{Reaction kinetics}

\begin{itemize}
  \item $rS\big(1-N/K\big)$ is logistic recruitment of susceptibles, with intrinsic growth rate $r$ and local carrying capacity $K(\mathbf{x})$. Density dependence acts through the \emph{total} local density $N$, so that infected and shedding animals compete for the same resources as susceptibles, but all newborns enter the susceptible class.
  \item $\beta_1 S R$ is \emph{direct} transmission, from contact between susceptible and shedding deer.
  \item $\beta_2 S W$ is \emph{indirect}, or \emph{environmental}, transmission, from susceptible deer encountering prion-contaminated ground. This term is the defining feature of CWD models relative to directly transmitted diseases: because the environmental reservoir persists independently of the host, it permits the disease to spread and to recur at deer densities far below the threshold required to sustain purely direct transmission.
  \item $\sigma I$ is progression from the sub-clinical, non-shedding class to the shedding class, with mean incubation period $1/\sigma$ --- for CWD characteristically long, many months to a few years.
  \item $\alpha R$ is removal of shedding deer through disease-induced mortality, together with any predation or harvest that disproportionately targets symptomatic animals. Removal is terminal in this model: there is no recovery pathway.
  \item $\rho R$ is deposition of prions into the environment, aggregating shedding during life with the contamination released at death and carcass decomposition.
  \item $\delta W$ is loss of environmental infectivity through prion degradation, burial, or transport out of the surface soil layer; $1/\delta$ is the effective environmental persistence time, which for CWD prions may be on the order of years.
\end{itemize}

The two transmission routes act on very different timescales. Direct transmission is mediated by the current shedding population, whereas $W$ integrates past shedding over a memory horizon of length $1/\delta$; the environment therefore acts as a slowly discharging store of infection, decoupling present risk from present prevalence. Prion contamination is treated as immobile, so \eqref{eq:W} is a purely local ordinary differential equation at each point of $\Gamma$, coupled to the deer compartments only through the pointwise value of $R$. This asymmetry shapes the numerical scheme of Section~\ref{sec:fem}, and Section~\ref{sec:adjoint-system} shows that the adjoint system inherits it exactly.

\subsection{Carrying capacity \texorpdfstring{$K$}{K}}
\label{sec:carrying}

The local deer carrying capacity is a base capacity $K_0$ (deer/km$^2$) modulated by a land-cover-dependent scaling coefficient $c(\mathbf{x})$:
\begin{equation}
  K(\mathbf{x}) = K_0\, c(\mathbf{x}),
  \label{eq:K}
\end{equation}
with $c$ taking one of the discrete values of Table~\ref{tab:landcover-values} according to the land cover class at $\mathbf{x}$. Forested and mixed edge habitat supports substantially higher deer densities than open agricultural or developed land, and the coefficients are calibrated against regional deer density estimates so that the population-weighted mean of $K$ reproduces observed densities across the study region. Since $c$ is piecewise constant, $K$ is discontinuous across class boundaries; this poses no difficulty for the weak formulation of Section~\ref{sec:fem}, in which $K$ enters only through a zeroth-order term.

\subsection{Terrain- and land-cover-aware diffusion tensor \texorpdfstring{$D$}{D}}
\label{sec:diffusion}

Deer movement, and hence the effective diffusivity of $S$, $I$ and $R$, depends on both land cover and local terrain steepness, and it is not isotropic: deer follow contours far more readily than they climb. We therefore treat diffusion along the local steepest-descent direction separately from diffusion along the perpendicular contour direction, so that the slope penalty applies to the former while the latter remains free, and scale both by the permeability of the local land cover.

Let $\mathbf{n}(\mathbf{x})$ be the unit normal to $\Gamma$, let $\mathbf{g}=(0,0,-1)$ be the downward vertical, and let $\theta$ be the terrain slope angle, defined by $\cos\theta = \mathbf{g}\cdot\mathbf{n}$, so that $|\cos\theta| = 1$ on flat ground and decreases as the terrain steepens. The local downhill tangent direction is the normalized projection of $\mathbf{g}$ onto the tangent plane,
\begin{equation}
  \mathbf{e}_s = \frac{\mathbf{g}-(\mathbf{g}\cdot\mathbf{n})\mathbf{n}}{\big\lVert \mathbf{g}-(\mathbf{g}\cdot\mathbf{n})\mathbf{n}\big\rVert},
\end{equation}
with a projected horizontal reference direction substituted wherever the surface is locally flat and $\mathbf{e}_s$ is numerically undefined, and $\mathbf{e}_c$ is the unit tangent perpendicular to it, so that $\mathbf{e}_s\otimes\mathbf{e}_s + \mathbf{e}_c\otimes\mathbf{e}_c = \mathbb{I}-\mathbf{n}\otimes\mathbf{n}$ is the projector onto the tangent plane. The diffusion tensor is
\begin{equation}
  D(\mathbf{x}) = \ell(\mathbf{x})\Big[D_s(\mathbf{x})\,\mathbf{e}_s\otimes\mathbf{e}_s + D_c(\mathbf{x})\,\mathbf{e}_c\otimes\mathbf{e}_c\Big].
  \label{eq:Dtens}
\end{equation}
Here $\ell(\mathbf{x})\in(0,1]$ is the dimensionless land-cover permeability of Table~\ref{tab:landcover-values}, $D_c$ the contour-direction diffusivity, and $D_s$ the downhill/uphill diffusivity. Both terms are outer products of a vector with itself, so $D = D^{\mathsf T}$ pointwise and the surface diffusion operator is self-adjoint --- a property used repeatedly in Sections~\ref{sec:adjoint-derivation} and~\ref{sec:discrete-adjoint}.

Along contours we take $D_c = \kappa$, constant, so that $\kappa$ carries the entire physical magnitude of the diffusivity and every other factor in \eqref{eq:Dtens} is dimensionless. Along the slope direction, deer avoid traversing steep terrain, for which the energetic cost of movement is high, and $D_s$ is suppressed accordingly through the smooth activation
\begin{equation}
  a(\theta) = \frac{1+\exp\!\big(k(\cos\theta_0-1)\big)}{1+\exp\!\big(k(\cos\theta_0-|\cos\theta|)\big)},
  \qquad
  D_s = \kappa\big[\,\iota + (1-\iota)\,a(\theta)\,\big].
  \label{eq:Ds}
\end{equation}
The numerator normalizes $a$ to unity at $|\cos\theta|=1$, so that $D_s = D_c = \kappa$ on flat ground and the tensor is isotropic there by construction; $a$ then decays sigmoidally, with $\theta_0$ centring the transition and $k$ setting its sharpness (Figure~\ref{fig:Ds}). The floor $\iota>0$ keeps $D_s$ bounded away from zero: without it the tangential part of \eqref{eq:Dtens} drops from rank two to rank one on the steepest facets, leaving the stiffness matrix singular along $\mathbf{e}_s$. It is thus a conditioning parameter rather than a movement parameter, a distinction taken up in Section~\ref{sec:params-movement}. Both $\theta_0$ and $k$ should be regarded as tunable and calibrated against telemetry-derived step-selection data where available.

Finally, each PDE compartment carries its own dimensionless mobility multiplier, $D_X = D_X^{\text{scale}}D(\mathbf{x})$ for $X\in\{S,I,R\}$, so that shedding deer may be modelled as less mobile than susceptibles.

\begin{figure}[h]
\centering
\includegraphics[width=3in]{Screenshot 2026-09-02 103330.png}
\caption{Slope-dependent function $a(\theta)$ of \eqref{eq:Ds} with $k=50$ and $\theta_0=16.3^\circ$}
\label{fig:Ds}
\end{figure}

\subsection{Parameters}

Table~\ref{tab:params} collects the model parameters, their interpretations, and their units; adopted numerical values and their provenance are given in Section~\ref{sec:params} and Table~\ref{tab:paramvalues}.

\begin{table}[h]
\centering
\caption{Model parameters. Deer densities are in deer/km$^2$, mesh lengths in metres (Section~\ref{sec:params-units}), and time in years.}
\label{tab:params}
\begin{tabular}{@{}llll@{}}
\toprule
Symbol & Description & Units & Value \\
\midrule
$r$        & Intrinsic growth rate of deer                 & yr$^{-1}$                        & Table~\ref{tab:paramvalues} \\
$K_0$      & Base carrying capacity                        & deer\,km$^{-2}$                  & Table~\ref{tab:paramvalues} \\
$c$        & Land-cover scaling of $K$                     & ---                              & Table~\ref{tab:landcover-values} \\
$\beta_1$  & Direct transmission rate                      & km$^{2}$\,deer$^{-1}$\,yr$^{-1}$ & Table~\ref{tab:paramvalues} \\
$\beta_2$  & Environmental transmission rate               & (units of $W$)$^{-1}$\,yr$^{-1}$ & Table~\ref{tab:paramvalues} \\
$\sigma$   & Progression rate, $I \to R$                   & yr$^{-1}$                        & Table~\ref{tab:paramvalues} \\
$\alpha$   & Removal rate of shedding deer                 & yr$^{-1}$                        & Table~\ref{tab:paramvalues} \\
$\rho$     & Environmental prion shedding rate             & (units of $W$)\,deer$^{-1}$\,yr$^{-1}$ & Table~\ref{tab:paramvalues} \\
$\delta$   & Environmental prion decay rate                & yr$^{-1}$                        & Table~\ref{tab:paramvalues} \\
$\kappa$   & Diffusivity scale (contour direction)         & m$^{2}$\,yr$^{-1}$               & Table~\ref{tab:paramvalues} \\
$\ell$     & Land-cover permeability                       & ---                              & Table~\ref{tab:landcover-values} \\
$\iota$    & Anisotropy (degeneracy) floor                 & ---                              & Table~\ref{tab:paramvalues} \\
$\theta_0,\,k$ & Slope-activation threshold and steepness  & ---                              & Table~\ref{tab:paramvalues} \\
$D_X^{\text{scale}}$ & Per-compartment mobility            & ---                              & Table~\ref{tab:paramvalues} \\
\bottomrule
\end{tabular}
\end{table}

\section{Parameter Values and Their Provenance}
\label{sec:params}

CWD dynamics unfold on much longer timescales than acute wildlife diseases: latency is measured in months to years and environmental persistence in years. Parameter values must therefore be drawn from the CWD-specific literature rather than carried over from models of acute, directly transmitted disease.

\subsection{Provenance classes}

Each value carries one of five tags, and every guessed value is shaded so it can be found at a glance.
\begin{itemize}
\item \tagLIT{} --- derived from published CWD field or modelling data, with the derivation shown.
\item \tagDEF{} --- a unit-defining normalization. Exact by construction, not an estimate.
\item \tagCAL{} --- back-solved from a stated calibration target (a target $R_0$, a target balance between terms in the objective). The \emph{target} is the assumption; the arithmetic is not.
\item \tagGUESS{} --- expert judgement with no supporting citation for this study region. These are the values to sweep.
\item \tagNUM{} --- a parameter that exists for conditioning or discretization and represents nothing about deer or prions; tuning one of these on ecological grounds is a category error.
\end{itemize}
Numerical values below were read from the parameter table of \citet{Almberg2011}, the one source whose tabulated parameters were consulted directly. \citet{Miller2006} and \citet{Wasserberg2009} are cited for their qualitative findings; reconciling their fitted transmission parameters and $R_0$ estimates against the values adopted here would move several \tagCAL{} entries toward \tagLIT{}.

\subsection{Units, and the two length scales in play}
\label{sec:params-units}

The mesh pipeline reprojects the DEM and the land-cover raster to a UTM coordinate reference system, so \textbf{mesh coordinates are in metres}, while deer densities are naturally quoted per square kilometre. The two conventions never meet: the transmission and growth terms are \emph{pointwise} and carry no area integral, so $S,I,R$ may be read as deer/km$^2$ and $\beta_1,\beta_2$ as km$^2$ per unit per year without any conversion, and the diffusion term is the only place where mesh length units enter, so $\kappa$ --- and $\kappa$ alone --- must be given in m$^2$/yr. The objective \eqref{eq:objective} integrates $d\Gamma$ in m$^2$, which multiplies $J$ by a constant and therefore leaves the minimizer unchanged.

\subsection{Host stage durations: $\sigma$ and $\alpha$}

\citet{Almberg2011} report, for white-tailed deer, a mean exposed (infected, non-shedding) period of 27 weeks (range 15--35), a mean infectious (shedding, pre-clinical) period of 36 weeks (range 25--44), and a mean clinical period of 17 weeks (range 1--36). This model carries only two infected classes, so $I$ is their exposed stage and $R$ lumps their infectious and clinical stages together. The lumping matters twice over: $1/\alpha$ is the full shedding duration rather than the much shorter clinical survival time, and the reduced mobility of clinically affected deer must be entered as a blend rather than as a clinical-stage value (Section~\ref{sec:params-movement}). It gives
\[
\sigma = \frac{52}{27} = 1.926\ \mathrm{yr}^{-1}
\qquad\text{and}\qquad
\alpha = \frac{52}{36+17} = \frac{52}{53} = 0.981\ \mathrm{yr}^{-1},
\]
i.e.\ a mean latency of 6.2 months and a mean shedding duration of 12.2 months. Propagating their reported ranges gives $\sigma \in [1.49, 3.47]$ and $\alpha \in [0.65, 1.95]$. Both are \tagLIT{}.

\subsection{No background mortality: what is omitted, and how to put it back}
\label{sec:params-mortality}

The model carries no non-CWD mortality term. $\sigma$ is the only exit from $I$ and $\alpha$ the only exit from $R$, so every infected deer reaches the shedding class with probability one, and $\alpha$ absorbs disease-induced death together with any predation or harvest that disproportionately targets symptomatic animals. This is a deliberate choice, and it has two consequences that shape how a result may be read. First, there is no harvest anywhere in the model, so the control $v$ of Section~\ref{sec:control} is \emph{total} removal effort on the shedding class rather than effort above a normal hunt, and $v\equiv0$ is an unmanaged herd rather than a status-quo managed one. Second, the optimizer sees every culled animal as a certain future shedder, which biases the optimal effort upward; the size of the effect is not small, since at the \citet{Almberg2011} weekly survival $s = 0.991$, i.e.\ $\mu = -52\ln(0.991) = 0.470\ \mathrm{yr}^{-1}$, roughly a fifth of infected animals would be removed before ever shedding.

Reinstating $\mu$ is a contained change, but worth spelling out because it is easy to do incorrectly. It means adding $-\mu I$ and $-\mu R$ to \eqref{eq:I}--\eqref{eq:R}, the matching $-\mu$ entries to the $(I,I)$ and $(R,R)$ diagonal of the Jacobian \eqref{eq:jacobian} and hence to the adjoint \eqref{eq:adjI}--\eqref{eq:adjR}, and --- the step that is easy to forget --- \emph{rescaling $\beta_1$ and $\beta_2$}, which were back-solved from a target $R_0$ whose formula contains $\mu$ (Section~\ref{sec:params-r0}). Bolting $\mu = 0.470$ onto the present coefficients without rescaling drops $R_0$ from $2.00$ to $1.09$, a whisker above threshold. One further point would need settling first: \citeauthor{Almberg2011}'s stage durations are used here as \emph{progression} times, with any survival term applied separately; if they are instead realized residence times that already net out mortality, then adding $\mu$ to $R$ would double-count and $\alpha$ would have to be reduced to compensate.

\subsection{Environmental persistence $\delta$, and the definition of a prion-unit}

\citet{Almberg2011} parameterize environmental prion loss by a weekly survival probability $\gamma$, which they vary over $[0.9481, 0.9983]$. Converting to a continuous decay rate, $\delta = -52\ln\gamma$, this spans
\[
\delta \in [0.089,\ 2.77]\ \mathrm{yr}^{-1},
\qquad\text{i.e.\ a persistence } 1/\delta \in [0.36,\ 11.3]\ \text{yr}.
\]
We adopt $\delta = 0.5\ \mathrm{yr}^{-1}$, a two-year persistence: the geometric midpoint of their explored range, which coincides with the direct experimental evidence of \citet{Miller2004}, who showed that mule deer paddocks remained infectious to na\"ive deer at least two years after the removal of infected animals. \tagLIT{}. This is the single most consequential parameter in the model --- precisely the one \citet{Almberg2011} identify as driving prevalence and population decline --- and it should be the first thing swept.

The shedding rate $\rho$ cannot be estimated independently, because the environmental compartment has no measurable absolute scale: multiplying $\rho$ and dividing $\beta_2$ by the same constant leaves the dynamics identical. We therefore fix the scale by \emph{definition},
\begin{equation}
\rho \equiv 1\ \text{prion-unit}/(\text{deer}\cdot\text{yr}),
\qquad\text{i.e.\ one prion-unit is the quantity shed by one shedding deer in one year,}
\end{equation}
\tagDEF{}, which absorbs the arbitrary scale and leaves $\beta_2$ as the single free environmental parameter. (\citeauthor{Almberg2011} instead carry a shedding rate of 2 units/week plus a 100-unit carcass pulse at death; those units are internal to their model and do not transfer.)

\subsection{Transmission coefficients $\beta_1$ and $\beta_2$, from a target $R_0$}
\label{sec:params-r0}

Published per-capita transmission coefficients for CWD are not directly transferable here. \citet{Almberg2011} report $\beta_d \in [7.5\times10^{-7}, 2.5\times10^{-5}]$ and $\beta_i \in [1.1\times10^{-8}, 1.75\times10^{-8}]$, but these are per-individual rates in a closed 20\,000-animal population with no spatial extent, and carry different dimensions from the density-dependent, per-area coefficients this model needs. We therefore back-solve $\beta_1$ and $\beta_2$ from a stated reproductive number, which is the quantity the two literatures actually share.

In a wholly susceptible landscape at $S=K_0$, an animal newly entering $I$ reaches the shedding class with certainty --- there is no competing risk (Section~\ref{sec:params-mortality}) --- and then remains there for a mean time $\Pi = 1/\alpha = 1.019$~yr, $\alpha$ being its only exit. During that time it infects susceptibles directly at rate $\beta_1 K_0$, and it deposits $\rho\,\Pi$ prion-units, each of which persists a mean $1/\delta$ years while infecting at rate $\beta_2 K_0$. Hence
\begin{equation}
R_0^{\mathrm{dir}} = \beta_1 K_0\,\Pi = \frac{\beta_1 K_0}{\alpha},
\qquad
R_0^{\mathrm{env}} = \frac{\beta_2 K_0\,\rho\,\Pi}{\delta} = \frac{\beta_2 K_0\,\rho}{\delta\,\alpha},
\qquad
R_0 = R_0^{\mathrm{dir}} + R_0^{\mathrm{env}}.
\label{eq:R0}
\end{equation}
Note the structure of $R_0^{\mathrm{env}}$: the environmental route's contribution is inflated by the factor $1/\delta$, the persistence time. This is the algebraic content of \citeauthor{Almberg2011}'s central observation that the effective infectious period can vastly exceed the host's lifespan, and hence that $R_0$ can be far larger than a direct-transmission analysis would suggest.

We target $R_0^{\mathrm{dir}} = R_0^{\mathrm{env}} = 1$, so that neither route alone is comfortably self-sustaining but the two together give $R_0 = 2$ --- the qualitative regime in which environmental persistence is what makes the epizootic go, and the regime of interest. A total of $R_0 = 2$ is a modest, deliberately conservative choice, at the low end of what the transmission-dynamics literature for cervid prion disease suggests \citep{Miller2006}. Inverting \eqref{eq:R0} with $\alpha = 0.981$, $\delta = 0.5$, $\rho = 1$ and $K_0 = 10$,
\[
\beta_1 = \frac{R_0^{\mathrm{dir}}\,\alpha}{K_0} = 0.0981\ \frac{\mathrm{km}^2}{\text{deer}\cdot\mathrm{yr}},
\qquad
\beta_2 = \frac{R_0^{\mathrm{env}}\,\delta\,\alpha}{K_0\,\rho} = 0.04905\ \frac{\mathrm{km}^2}{\text{prion-unit}\cdot\mathrm{yr}}.
\]
Both are \tagCAL{}. \emph{The even 1:1 split is the assumption, not the numbers} --- it is a free choice and should be swept. Note also that $\beta_1$ and $\beta_2$ are tied to $K_0$ and $\alpha$, and $\beta_2$ also to $\delta$ and $\rho$: revise any of these and both coefficients must be rescaled to hold $R_0$ fixed, or the epidemic threshold moves as a side effect of an edit that had nothing to do with transmission. Section~\ref{sec:params-mortality} describes exactly that failure mode.

At these values the spatially uniform endemic equilibrium is $S^* = K_0/R_0 = 5.0$, $I^* = 0.358$, $R^* = 0.703$, giving $N^* = 6.06$ deer/km$^2$ --- a $39\%$ population decline at $17.5\%$ prevalence, which is in the observed range for a CWD core area. That $S^* = K_0/R_0$ exactly is the usual threshold identity, and a useful check on any re-calibration.

\subsection{Demography: $r$ and $K_0$}

For $r$ we take $0.35\ \mathrm{yr}^{-1}$, a finite rate of increase $\lambda = e^{0.35} = 1.42$ for an unexploited herd well below carrying capacity. This is consistent with the per-capita birth rate of $0.55\ \mathrm{yr}^{-1}$ used by \citet{Almberg2011} net of unhunted adult mortality of order $0.15$--$0.20\ \mathrm{yr}^{-1}$; a plausible range is $0.25$--$0.50$. \tagLIT{} (derived, moderate confidence).

For $K_0$ we take $10$ deer/km$^2$ ($\approx 26$ deer/mi$^2$) in the best habitat class, which after applying Table~\ref{tab:landcover-values} gives a landscape mean nearer $6$--$8$ deer/km$^2$ --- a high but attainable density for a Midwestern or eastern white-tailed deer herd. \tagGUESS{}: it is region-specific and no source has been consulted for the actual study area, so it should be replaced with the managing agency's overwinter density estimate. Recall from Section~\ref{sec:params-r0} that changing $K_0$ obliges a matching rescale of $\beta_1$ and $\beta_2$.

\subsection{Movement: $\kappa$, the anisotropy floor, and compartment mobilities}
\label{sec:params-movement}

No fitted diffusion coefficient for white-tailed deer is available, so $\kappa$ is bracketed by two independent order-of-magnitude arguments.
\begin{enumerate}
\item \emph{From dispersal.} Roughly a quarter of white-tailed deer disperse once, as yearlings, over a median distance of order 8--10~km. Spreading that mean squared displacement over a population turnover time of a few years gives $D$ of order $3$~km$^2$/yr.
\item \emph{From front speed.} A Fisher--KPP invasion advances at $c = 2\sqrt{D\,r_I}$, where $r_I$ is the initial exponential growth rate of the infected class. An observed CWD front speed of a few km/yr with $r_I$ of order $0.3\ \mathrm{yr}^{-1}$ implies $D$ of order $5$--$8$~km$^2$/yr.
\end{enumerate}
Both land in a $2$--$8$~km$^2$/yr band, and we adopt $\kappa = 4\times10^{6}\ \mathrm{m}^2/\mathrm{yr} = 4$~km$^2$/yr on flat, maximally permeable land cover. \tagGUESS{} --- this is the least defensible parameter in the model. Its saving grace is that the spread rate goes as $\sqrt{\kappa}$, so a factor-of-four error in $\kappa$ is only a factor-of-two error in front speed.

The per-compartment multipliers $D_X^{\text{scale}}$ are dimensionless mobility ratios. We set the susceptible and infected scales to $1.0$ --- sub-clinical infected deer are not detectably less mobile than susceptibles --- and the $R$ scale to $0.3$. \tagGUESS{}. Because $R$ lumps pre-clinical shedding with the clinical phase, $0.3$ is a blend of near-normal and near-immobile rather than a clinical-stage value; a genuinely clinical-stage mobility would be far lower, and recovering it would require splitting $R$ in two.

The floor $\iota$ is retained at $0.02$, and it is worth being explicit that \textbf{$\iota$ is not a movement parameter and must not be tuned as one.} All of the slope physics lives in the activation function $a(\theta)$ of \eqref{eq:Ds}; what $\iota$ does is prevent $D_s$ from reaching zero and the tensor from degenerating. It is wanted small enough not to interfere with $a(\theta)$ and large enough to keep the tensor conditioned: $\iota = 0.02$ gives a tangent-plane condition number of $50$. It also sets the floor on mesh edge length, $h_{\min} = \iota\,h_{\max} = 3$~m at the adopted $h_{\max} = 150$~m, since the same activation drives the \texttt{mmg} refinement metric. \tagNUM{}. Reading this regularization as a physical anisotropy ratio is an easy mistake to make, since the two are multiplied together in the same expression.

The activation threshold $\cos\theta_0 = 0.96$ centres the transition at $\theta_0 = 16.3^\circ$, inside the 15--25$^\circ$ band above which cervid habitat-selection studies generally show avoidance, and the steepness $k = 50$ spreads that transition over slopes of $0$--$23^\circ$. Both \tagGUESS{}; unlike $\iota$, these two \emph{are} the movement parameters, and $k$ in particular controls only how abruptly mobility falls off, not where.

\subsection{The land-cover tables}
\label{sec:params-landcover}

Table~\ref{tab:landcover-values} gives the adopted land-cover coefficients. \emph{The two columns measure different things and are not to be collapsed into one.} The coefficient $c$ measures \textbf{habitat quality} --- how many deer a cover type supports, entering $K(\mathbf{x}) = K_0 c(\mathbf{x})$ --- while $\ell$ measures \textbf{permeability}, how readily deer move through it, entering the diffusion tensor \eqref{eq:Dtens}. For several classes the two are anti-correlated. Open grassland is highly permeable but poor white-tailed deer habitat: it offers little browse and no cover, so deer cross it quickly rather than settling in it. Woody wetlands are the reverse --- prime cover and winter browse, but slow going. Developed open space (the exurban lawn-and-woodlot mosaic, parks, golf courses) supports famously high suburban deer densities while being only moderately permeable.

Every entry of Table~\ref{tab:landcover-values} is a \tagGUESS{}, and both columns should be calibrated against regional density estimates stratified by cover type. Two entries deserve comment on numerical rather than ecological grounds. Open water is given permeability $10^{-2}$ rather than a hard zero: deer do swim, so large water bodies are strong but not absolute barriers, and a hard zero would leave the diffusion tensor \eqref{eq:Dtens} singular over every water cell. And both defaults, applied to class codes absent from the table, are $0.5$, so that unmapped or nodata cells fall back to a mid-range value rather than to prime, maximally permeable habitat.

\begin{table}[h]
\centering
\small
\caption{Land cover coefficients, keyed to NLCD/MRLC class codes. Column $c$ is the carrying-capacity scaling in $K(\mathbf{x})=K_0\,c(\mathbf{x})$ of Section~\ref{sec:carrying}; column $\ell$ is the dimensionless permeability multiplier in \eqref{eq:Dtens}, normalized so the most permeable classes sit at $1.0$. Every value in this table is a \tagGUESS{}, hence the shading.}
\label{tab:landcover-values}
\rowcolors{2}{guessbg}{guessbg}
\begin{tabular}{@{}llcc@{}}
\toprule
Code & Land cover class & $c$ (habitat) & $\ell$ (permeability) \\
\midrule
0  & Unknown or nodata            & 0.50 & 0.50 \\
11 & Open water                   & 0.00 & 0.01 \\
12 & Perennial ice or snow        & 0.00 & 0.30 \\
21 & Developed, open space        & 0.90 & 0.60 \\
22 & Developed, low intensity     & 0.60 & 0.40 \\
23 & Developed, medium intensity  & 0.25 & 0.15 \\
24 & Developed, high intensity    & 0.05 & 0.02 \\
31 & Barren land                  & 0.10 & 0.90 \\
41 & Deciduous forest             & 1.00 & 0.70 \\
42 & Evergreen forest             & 0.60 & 0.60 \\
43 & Mixed forest                 & 0.95 & 0.65 \\
52 & Shrub or scrub               & 0.85 & 0.80 \\
71 & Grassland or herbaceous      & 0.50 & 1.00 \\
81 & Pasture or hay               & 0.70 & 1.00 \\
82 & Cultivated crops             & 0.75 & 0.90 \\
90 & Woody wetlands               & 0.90 & 0.50 \\
95 & Emergent herbaceous wetlands & 0.35 & 0.35 \\
\bottomrule
\end{tabular}
\end{table}

\subsection{Horizon and time step}

The planning horizon is $T = 20$~yr. \tagGUESS{} --- a management choice rather than a measured quantity, but CWD epizootics take one to three decades to move from introduction to double-digit prevalence, so 20 years is about the shortest window over which a culling policy can be meaningfully evaluated. The time step is $\Delta t = 0.02$~yr $= 7.3$~days, matching the weekly step of \citet{Almberg2011}. \tagLIT{}. Accuracy, not stability, is the binding constraint: diffusion is implicit and the $W$ update is exact, so the only explicit rates are the reaction terms, the fastest of which is $\sigma$, giving $\Delta t\,\sigma = 0.039 \ll 1$. Both memory and runtime scale with $T/\Delta t$, here 1000 steps, since the control driver stores the entire state trajectory.

\subsection{Initial conditions}
\label{sec:params-ic}

The susceptible field is \emph{not} initialized at the carrying capacity. $K(\mathbf{x})$ is a piecewise-constant land-cover lookup and is therefore not a steady state of the $S$ equation: diffusion smooths it across every land-cover patch boundary, and over water, where $K=0$ and recruitment is replaced by decay at rate $-r$, the landscape acts as an absorbing sink that draws the profile down for a diffusion length around every shoreline. Starting at $S=K$ superimposes an order-one \emph{demographic} transient on the epidemic one, and the two are not separated in time: logistic relaxation has an $e$-folding time $1/r = 2.9$~yr, against an epizootic that develops over one to three decades, so the transient is still running while the quantities of interest are being measured. For the control problem of Section~\ref{sec:control} this is worse than cosmetic, since the spurious structure it creates is something the optimizer will spend control effort on.

We instead take $S(\cdot,0)=S_\infty$, the \emph{disease-free equilibrium}: the solution of
\begin{equation}
\Div\big(D_S\,\Grad S_\infty\big) + rS_\infty\Big(1-\frac{S_\infty}{K}\Big) = 0
\label{eq:disease-free}
\end{equation}
on the habitat portion of $\Gamma$, with the growth term replaced by $-rS_\infty$ wherever $K<0.01$ and under the same zero-flux boundary condition as the full system, obtained numerically by relaxing the $S$ equation alone from $S=K$ (Section~\ref{sec:spinup}). Two numerical parameters govern that relaxation. The stopping tolerance $\varepsilon_{\text{drift}} = 10^{-5}$~yr$^{-1}$ bounds $\max_\Gamma|\partial_t S|$ relative to $\max_\Gamma K$, so at $K_0 = 10$~deer/km$^2$ the field is required to be moving by less than $10^{-4}$~deer/km$^2$/yr --- negligible against every other error in the model. The cap $T_{\text{spin}} = 60$~yr is a ceiling, not a target: the relaxation exits as soon as the tolerance is met. It is generous because the slow process is diffusive rather than logistic, equilibration across a habitat patch of linear scale $L$ taking of order $L^2/\kappa$ --- 1~yr for a 2~km patch but 25~yr for a 10~km one, with shoreline drawdown around a large water body slower still, which is why the stopping rule is a convergence test rather than a fixed duration. Both are \tagNUM{}.

The infected class is initialized as a superposition of isotropic Gaussian outbreak foci,
\[
I(\mathbf{x},0) = \sum_{g} m_g \exp\!\Big(-\frac{\lVert \mathbf{x}_{\parallel}-\mathbf{c}_g\rVert^2}{2\varsigma_g^2}\Big),
\]
evaluated in horizontal coordinates, with $R$ and $W$ initialized to zero everywhere. Each focus has standard deviation $\varsigma = 500$~m --- roughly the linear scale of a white-tailed deer home range of 1--3~km$^2$, i.e.\ the area over which a single index case actually makes contacts and deposits prions (\tagGUESS{}) --- and an amplitude fixed by requiring the focus to integrate to \emph{exactly one} introduced animal: since $\int m\,e^{-\lVert x\rVert^2/2\varsigma^2}\,dA = 2\pi\varsigma^2 m$ and $\varsigma = 0.5$~km,
\[
m = \frac{1}{2\pi(0.5)^2} = 0.637\ \text{deer/km}^2. \qquad \tagCAL{}
\]
The focus \emph{locations} are arbitrary and mesh-dependent (\tagGUESS{}); they should be replaced by actual first-detection coordinates for the study region.

\subsection{Summary table}

\begin{table}[h]
\centering
\small
\caption{Model parameters, adopted values, and provenance. Shaded rows are values with no supporting citation for this study region --- the ones to sweep. Cost weights and control bounds are discussed in Section~\ref{sec:control}.}
\label{tab:paramvalues}
\begin{tabular}{@{}llllp{4.6cm}@{}}
\toprule
Symbol & Description & Value & Units & Provenance \\
\midrule
$\sigma$ & Latency $\to$ shedding rate & $1.926$ & yr$^{-1}$ & \tagLIT{} $52/27$ wk, \citet{Almberg2011}; range $1.49$--$3.47$ \\
$\alpha$ & CWD removal rate of shedding deer & $0.981$ & yr$^{-1}$ & \tagLIT{} $52/53$ wk, \citet{Almberg2011}; range $0.65$--$1.95$ \\
$\delta$ & Environmental prion decay & $0.5$ & yr$^{-1}$ & \tagLIT{} $-52\ln\gamma$, \citet{Almberg2011}; corroborated by \citet{Miller2004}; range $0.089$--$2.77$ \\
$\rho$ & Prion shedding rate & $1.0$ & PU/(deer$\cdot$yr) & \tagDEF{} defines the prion-unit; exact \\
$\beta_1$ & Direct transmission & $0.0981$ & km$^2$/(deer$\cdot$yr) & \tagCAL{} from $R_0^{\mathrm{dir}}=1$, Eq.~\eqref{eq:R0} \\
$\beta_2$ & Environmental transmission & $0.04905$ & km$^2$/(PU$\cdot$yr) & \tagCAL{} from $R_0^{\mathrm{env}}=1$, Eq.~\eqref{eq:R0} \\
$r$ & Intrinsic growth rate & $0.35$ & yr$^{-1}$ & \tagLIT{} derived from \citet{Almberg2011} fecundity; range $0.25$--$0.50$ \\
$\Delta t$ & Time step & $0.02$ & yr & \tagLIT{} weekly step of \citet{Almberg2011} \\
$\varepsilon_{\text{drift}}$ & Spin-up stopping tolerance & $10^{-5}$ & yr$^{-1}$ & \tagNUM{} $\max|\partial_t S|$ relative to $\max K$; Section~\ref{sec:spinup} \\
$T_{\text{spin}}$ & Spin-up duration cap & $60$ & yr & \tagNUM{} a ceiling, not a target; exits on $\varepsilon_{\text{drift}}$ \\
\guessrow $K_0$ & Base carrying capacity & $10$ & deer/km$^2$ & \tagGUESS{} region-specific; replace with agency estimate \\
\guessrow $\kappa$ & Diffusivity scale & $4\times10^{6}$ & m$^2$/yr & \tagGUESS{} order of magnitude from dispersal and front speed \\
$\iota$ & Tensor degeneracy floor & $0.02$ & --- & \tagNUM{} regularization, \emph{not} a movement parameter; also sets $h_{\min}$ \\
\guessrow $\cos\theta_0$ & Activation threshold & $0.96$ & --- & \tagGUESS{} centres the transition at $16.3^\circ$ \\
\guessrow $k$ & Activation steepness & $50$ & --- & \tagGUESS{} smoothness only \\
\guessrow $D_{S,I,R}^{\text{scale}}$ & Compartment mobilities & $1,\,1,\,0.3$ & --- & \tagGUESS{} $R$ value is a lumped-stage blend \\
\guessrow $c(\mathbf{x}),\ \ell(\mathbf{x})$ & Land-cover coefficients & Table~\ref{tab:landcover-values} & --- & \tagGUESS{} expert judgement throughout \\
\guessrow $T$ & Planning horizon & $20$ & yr & \tagGUESS{} management choice \\
\guessrow $c_1,c_2,c_3$ & Objective weights & $1,\,0.5,\,2$ & --- & \tagGUESS{} term-balance recipe, Section~\ref{sec:control}; $c_1$ weights $R$ \\
\guessrow $c_4$ & Linear (per-deer) control weight & $0$ & --- & \tagGUESS{} sets the effort threshold in \eqref{eq:pointwise-optimality}; $0$ by default \\
\guessrow $v_{\max}$ & Maximum culling effort on $R$ & $1.0$ & yr$^{-1}$ & \tagGUESS{} $63\%$ annual removal at full effort \\
\bottomrule
\end{tabular}
\end{table}

\subsection{What is still a guess}

Reading the shading of Table~\ref{tab:paramvalues}: the \emph{disease} parameters are grounded --- the stage durations, the environmental decay rate and the host growth rate all trace to published CWD data, and the two transmission coefficients follow from those by an explicit $R_0$ calculation whose only free input is the direct/environmental split. What remains unsupported is everything \emph{spatial} and everything \emph{economic}: the carrying capacity, the diffusivity scale, the two slope-activation parameters, the entire land-cover table, and every objective and control-bound parameter. (The three \tagNUM{} rows are unshaded because they are not part of this argument: no field measurement would settle any of them.) That shapes what conclusions the model can currently support. Results about the \emph{timing} and \emph{qualitative structure} of an optimal culling policy rest on the grounded parameters and are reasonably robust; results about the absolute spatial \emph{extent} of spread, or about the cost-effectiveness of a policy in any real currency, are not.

Sensitivity work should therefore begin with $\delta$ --- the parameter \citet{Almberg2011} identify as controlling the outcome, and the one with the widest published range --- followed by the $R_0^{\mathrm{dir}}:R_0^{\mathrm{env}}$ split, then $\kappa$, then $K_0$ with $\beta_1,\beta_2$ rescaled to hold $R_0$ fixed.

\section{Numerical Methods: Surface Finite Element Discretization}
\label{sec:fem}

We solve \eqref{eq:S}--\eqref{eq:W} using a surface finite element method on a triangulated approximation $\Gamma_h$ of $\Gamma$, built directly from the DEM of Section~\ref{sec:domain}: elevation samples are triangulated in the horizontal plane and lifted to $\mathbb{R}^3$, producing a mesh of flat triangles whose union approximates $\Gamma$. Each vertex inherits an interpolated land cover class, used to evaluate $K$ and $\ell$, and a discrete unit normal, computed as an area-weighted average of the normals of incident triangles and used to evaluate the diffusion tensor \eqref{eq:Dtens}.

\subsection{Weak formulation}
\label{sec:weak}

Let $V_h \subset H^1(\Gamma_h)$ be the space of continuous, piecewise-linear finite element functions on $\Gamma_h$, with nodal basis $\{\varphi_i\}_{i=1}^N$. Multiplying \eqref{eq:S}--\eqref{eq:R} by a test function $\psi \in V_h$ (we reserve $v$ for the culling control of Section~\ref{sec:control}), integrating over $\Gamma_h$, and applying the surface divergence theorem to the diffusion term gives: find $S,I,R \in V_h$ such that for all $\psi \in V_h$,
\begin{align}
\int_{\Gamma_h} \dot S\, \psi\, dA &= -\int_{\Gamma_h} D_S\,\Grad S \cdot \Grad \psi \, dA + \int_{\Gamma_h}\Big[rS\big(1-\tfrac{N}{K}\big) - \beta_1 SR - \beta_2 SW\Big] \psi \, dA, \\
\int_{\Gamma_h} \dot I\, \psi\, dA &= -\int_{\Gamma_h} D_I\,\Grad I \cdot \Grad \psi \, dA + \int_{\Gamma_h}\big[\beta_1 SR + \beta_2 SW - \sigma I\big] \psi \, dA, \\
\int_{\Gamma_h} \dot R\, \psi\, dA &= -\int_{\Gamma_h} D_R\,\Grad R \cdot \Grad \psi \, dA + \int_{\Gamma_h}\big[\sigma I - \alpha R\big] \psi \, dA,
\end{align}
the zero-flux conditions \eqref{eq:bc} being imposed naturally by the omission of boundary terms. Equation~\eqref{eq:W} requires no weak formulation: it is discretized directly at the nodes,
\begin{equation}
\dot W_i = \rho R_i - \delta W_i, \qquad i=1,\dots,N,
\end{equation}
so that the $W$ update is purely local.

\subsection{Semi-discrete system}

Expanding $S_h = \sum_i S_i(t)\varphi_i$, and similarly for $I_h$ and $R_h$, yields
\begin{align}
M\dot{\mathbf{S}} &= -A_S\,\mathbf{S} + \mathbf{F}_S(\mathbf{S},\mathbf{I},\mathbf{R},\mathbf{W}), \qquad
&&M\dot{\mathbf{R}} = -A_R\,\mathbf{R} + \mathbf{F}_R(\mathbf{I},\mathbf{R}), \\
M\dot{\mathbf{I}} &= -A_I\,\mathbf{I} + \mathbf{F}_I(\mathbf{S},\mathbf{I},\mathbf{R},\mathbf{W}), \qquad
&&\dot{\mathbf{W}} = \rho\,\mathbf{R} - \delta\,\mathbf{W},
\end{align}
where $M_{ij} = \int_{\Gamma_h}\varphi_i\varphi_j\,dA$ is the consistent mass matrix, $[A_X]_{ij} = \int_{\Gamma_h} D_X\,\Grad\varphi_i\cdot\Grad\varphi_j\,dA$ the stiffness matrix built with the tensor \eqref{eq:Dtens}, and $\mathbf{F}_S,\mathbf{F}_I,\mathbf{F}_R$ the $L^2$-projected reaction terms.

In place of $M$ we use the lumped mass matrix $M_L = \mathrm{diag}(M\mathbf{1})$, obtained by integrating the mass terms with a vertex quadrature rule: at the vertices $\varphi_i\varphi_j = \delta_{ij}$, so the off-diagonal entries vanish while the row sums, and hence the total mass, are preserved exactly. Lumping makes the mass matrix trivially invertible and --- more importantly --- preserves positivity. The consistent mass matrix carries strictly positive off-diagonal entries, which destroy the M-matrix structure of $M+\Delta t A_X$ and allow the implicit diffusion step to undershoot on sharp initial data; the resulting negative densities reverse the sign of the logistic term and drive the solution to blow up. With $M_L$ the off-diagonals are zero and, on a mesh with no obtuse angles, $M_L + \Delta t A_X$ is an M-matrix, so a non-negative right-hand side yields a non-negative solution. For piecewise-linear elements lumping is an $O(h^2)$-consistent quadrature and therefore costs no order of accuracy. Both sides of the update are lumped; the nonlinear reaction terms are retained on a Gauss rule rather than collapsed to nodal evaluations.

Both $M$ and $A_X$ are symmetric, the latter because $D=D^{\mathsf T}$ pointwise (Section~\ref{sec:diffusion}). This is not merely an efficiency remark: Section~\ref{sec:discrete-adjoint} shows that it is exactly what allows the adjoint sweep to reuse the forward operator's factorization.

\subsection{Time integration}

The diffusion terms in $S,I,R$ are linear and potentially stiff, while the reaction and coupling terms are nonlinear but non-stiff, so we integrate with an implicit--explicit (IMEX) scheme. A first-order (semi-implicit Euler) step from $t^n$ to $t^{n+1}=t^n+\Delta t$ reads
\begin{equation}
\big(M_L + \Delta t\,A_S\big)\,\mathbf{S}^{n+1} = M_L\,\mathbf{S}^n + \Delta t\,\mathbf{F}_S(\mathbf{S}^n,\mathbf{I}^n,\mathbf{R}^n,\mathbf{W}^n),
\label{eq:imex}
\end{equation}
and analogously for $\mathbf{I}^{n+1}$ and $\mathbf{R}^{n+1}$. Having no diffusion, the environmental compartment is advanced by the exact solution of its linear decay-plus-source ODE at each node,
\begin{equation}
W_i^{n+1} = W_i^n e^{-\delta \Delta t} + \rho R_i^n\,\frac{1-e^{-\delta\Delta t}}{\delta},
\end{equation}
which is unconditionally stable regardless of $\Delta t$. Higher-order IMEX Runge--Kutta schemes may be substituted for $S,I,R$ to relax the time step restriction while retaining implicit diffusion; the exact $W$ update remains valid at any order.

\section{Implementation}
\label{sec:implementation}

The model and its discretization are implemented in \textsc{FEniCSx} (\texttt{DOLFINx}), using \texttt{basix.ufl} for element and mixed-space construction, \texttt{UFL} for the variational forms, and \texttt{PETSc} for linear algebra. Meshes are imported from Gmsh and solution fields written to XDMF/HDF5 for visualization in ParaView. Land cover classes and the derived fields $\ell(\mathbf{x})$ and $K(\mathbf{x})$ are precomputed by aligning the MRLC raster to mesh vertex coordinates (Section~\ref{sec:domain}) and supplied as one value per mesh vertex. Because those arrays are indexed by local degree of freedom and carry no global node identifiers, they are meaningful only for the partition that wrote them, and the solver therefore runs in serial; lifting the restriction means giving the mesh bundle a partition-independent index, not changing the discretization.

\subsection{Notation correspondence}

Table~\ref{tab:code-notation} maps the symbols of Section~\ref{sec:model} onto the identifiers used in the source, which differ in a few places: the shedding compartment $R$ is named \texttt{D} (``dying'') and the environmental compartment $W$ is named \texttt{E} (``environment''). One correspondence is a genuine trap when reading the code. \texttt{rho} is the \emph{environmental transmission} rate --- the role of $\beta_2$ here --- while the \emph{shedding} rate $\rho$ of Section~\ref{sec:model} is named \texttt{p\_env}, the opposite of the Greek-letter assignment used in the text. Table~\ref{tab:code-notation} is the authority.

\begin{table}[h]
\centering
\small
\caption{Correspondence between the mathematical symbols of Sections~\ref{sec:model}--\ref{sec:params} and identifiers in the DOLFINx implementation and its parameter file.}
\label{tab:code-notation}
\begin{tabular}{@{}lll@{}}
\toprule
Math & Code identifier & Role \\
\midrule
$S,I,R,W$ & \texttt{uS, uI, uD, uE} (subfunctions of \texttt{U}) & state variables \\
$r$ & \texttt{r} & intrinsic growth rate \\
$K(\mathbf{x})$ & \texttt{K\_func} & carrying capacity field, Eq.~\eqref{eq:K} \\
$c(\mathbf{x})$ & \texttt{land\_cover\_to\_carrying\_capacity} & class-to-capacity lookup, Table~\ref{tab:landcover-values} \\
$\ell(\mathbf{x})$ & \texttt{lcD\_scale} & land-cover permeability, Eq.~\eqref{eq:Dtens} \\
$S_\infty(\mathbf{x})$ & \texttt{S\_equilibrium} & disease-free susceptible field, Section~\ref{sec:spinup} \\
$\beta_1$ & \texttt{beta} & direct transmission rate \\
$\beta_2$ & \texttt{rho} & environmental transmission rate \\
$\sigma$ & \texttt{sigma} & latent $\to$ shedding progression rate \\
$\alpha$ & \texttt{alpha} & CWD removal rate of shedding deer \\
$\rho$ & \texttt{p\_env} & environmental prion shedding rate \\
$\delta$ & \texttt{delta\_e} & environmental prion decay rate \\
$D(\mathbf{x})$ & \texttt{DTens} & diffusion tensor \eqref{eq:Dtens}, Section~\ref{sec:diffusion} \\
$D_{S,I,R}^{\text{scale}}$ & \texttt{DS, DI, DD} (\texttt{compartment\_scales}) & per-compartment mobilities \\
$\kappa$ & \texttt{kappa} & diffusivity scale \\
$\iota$ & \texttt{isotropy} & anisotropy (degeneracy) floor \\
$\cos\theta_0$ & \texttt{activation\_cosine\_threshold} & slope-activation threshold \\
$k$ & \texttt{activation\_steepness} & slope-activation steepness \\
$T,\ \Delta t$ & \texttt{final\_time\_years}, \texttt{time\_step\_years} & horizon and time step \\
\bottomrule
\end{tabular}
\end{table}

\subsection{Monolithic assembly and block factorization}
\label{sec:monolithic}

All four compartments are carried in a single mixed space $P = V_S \times V_I \times V_R \times V_W$ built from four degree-1 Lagrange elements. We write $U = (u_S,u_I,u_R,u_W) \in P$ for the mixed trial function and $\Psi = (\psi_S,\psi_I,\psi_R,\psi_W) \in P$ for the mixed test function, capitals denoting the mixed functions and lower case their scalar components. The implicit side of the IMEX step \eqref{eq:imex} is then the bilinear form
\begin{equation}
a(U,\Psi) = \sum_{X\in\{S,I,R\}} \int_{\Gamma_h}\Big[\,u_X \psi_X + \Delta t\,(D_X\,\Grad u_X)\cdot\Grad \psi_X\Big]\,dA \;+\; \int_{\Gamma_h} u_W \psi_W\, dA .
\label{eq:bilinear}
\end{equation}
Being bilinear, $a$ is of course a function of $U$; what matters for cost is that its \emph{coefficients} are not functions of the state. Only the mesh geometry, the land cover and the diffusion parameters enter them, so testing \eqref{eq:bilinear} over the nodal basis of $P$ produces the same matrix at every time step. That matrix, written $B$, is therefore assembled and factorized \emph{once} before the time loop and reused unchanged, while the load vector on the right --- which does depend on the previous state --- is rebuilt each step. Mass terms are integrated with the vertex rule of Section~\ref{sec:fem}, giving $M_L$; stiffness terms use a Gauss rule.

Since \eqref{eq:bilinear} contains no cross-compartment terms, $B$ is block diagonal,
\begin{equation}
B = \operatorname{blockdiag}\big(\,\underbrace{M_L + \Delta t\,A_S}_{B_S},\;\underbrace{M_L + \Delta t\,A_I}_{B_I},\;\underbrace{M_L + \Delta t\,A_R}_{B_R},\;\underbrace{M_L}_{B_W}\,\big),
\label{eq:blockdiag}
\end{equation}
and is never assembled as a single matrix. Factorizing the blocks separately on the scalar space is strictly cheaper: $B_W$ is diagonal, so assembling monolithically would push a quarter of the unknowns through a sparse triangular solve to accomplish a vector division; blocks with equal mobility are the same matrix, so at the adopted $D_S^{\text{scale}} = D_I^{\text{scale}} = 1$ only two distinct factorizations exist; and sparse-factorization fill grows superlinearly. Solving blockwise is an exact restatement of the same linear system, each block is symmetric, and hence so is $B$.

The linear form collects all reaction and coupling terms evaluated explicitly at the previous time level, matching \eqref{eq:imex}, with the culling term $-vR$ added to the $R$-row in the controlled solver of Section~\ref{sec:control-implementation}. One detail matters downstream: a conditional switches the $S$-row growth term to pure exponential decay $-rS$ wherever $K(\mathbf{x})<0.01$, avoiding division by a near-zero carrying capacity over water and other non-habitat cells. Only the load vector is reassembled at each step, so a time step costs one vector assembly plus one back-substitution against the cached factors.

\subsection{Disease-free spin-up}
\label{sec:spinup}

The susceptible initial condition $S_\infty$ of Section~\ref{sec:params-ic} is not available in closed form --- \eqref{eq:disease-free} is a nonlinear elliptic problem on the same anisotropic tensor as the full system --- so it is computed by relaxation before the time loop. With $I=R=W\equiv0$ the mixed system collapses to its $S$ block, and the same lumped-mass IMEX step applies unchanged:
\begin{equation}
\big(M_L + \Delta t\,A_S\big)\,\mathbf{S}^{n+1} \;=\; M_L\mathbf{S}^{n} \;+\; \Delta t\,\mathbf{f}(\mathbf{S}^n),
\label{eq:spinup-step}
\end{equation}
where $\mathbf{f}$ assembles the same conditional growth term as the $S$-row above. The iteration starts from $\mathbf{S}^0 = \mathbf{K}$ and runs until
\begin{equation}
\frac{\max_i\big|S_i^{n+1}-S_i^{n}\big|}{\Delta t\,\max_i K_i} \;<\; \varepsilon_{\text{drift}},
\label{eq:spinup-stop}
\end{equation}
or until the cap $T_{\text{spin}}$ is reached, in which case the solver warns that the field is not converged. Normalizing by $\max_i K_i$ makes the criterion relative, so the tolerance carries no density units and need not be re-tuned when $K_0$ changes.

Mass lumping matters more here than anywhere else. The iteration starts from $\mathbf{K}$, which is piecewise constant and therefore as sharp as data on this mesh can be; with the consistent mass matrix the first implicit step would undershoot at the patch boundaries, and a negative $S$ flips the sign of the logistic term and diverges.

Because $S_\infty$ depends only on the mesh, the two land-cover arrays and the parameters $r,\kappa,\iota,\theta_0,k,D_S^{\text{scale}},\Delta t,\varepsilon_{\text{drift}}$ --- and on nothing that varies between runs, the control included --- it is computed once per mesh bundle and cached beside the mesh, with a companion record of exactly those inputs. A later run reuses the cached field only if that record still matches, so that editing the land cover or $\kappa$ cannot silently leave a stale equilibrium in place.

% ======================================================================
\section{An Optimal Control Problem: Targeted Culling}
\label{sec:control}
% ======================================================================

The model of Sections~\ref{sec:model}--\ref{sec:implementation} predicts what happens if nothing is done. We now use it as a constraint in a management problem: given a finite planning horizon $[0,T]$ and a culling operation whose intensity may be varied across the landscape and over time, where and when should that effort be spent?

\subsection{The control and the controlled state system}

We introduce a \emph{culling control} $v : \Gamma\times(0,T) \to \mathbb{R}$, a per-capita removal rate (yr$^{-1}$) applied to the \emph{shedding} class $R$. Culling therefore enters as an additional loss term $-vR$ in \eqref{eq:R}, removing shedding animals from the population entirely rather than moving them to another compartment. Since the model carries no harvest term of its own (Section~\ref{sec:params-mortality}), $v$ is \emph{total} removal effort on the shedding class, and $v\equiv 0$ is an unmanaged herd.

Targeting $R$ rather than $I$ is the choice that makes the control problem worth posing. $R$ is the only compartment that drives \emph{both} transmission routes --- directly through $\beta_1 SR$ and indirectly by generating the environmental reservoir through $\rho R$ --- so removing a shedding animal cuts both at once and, through $W$, keeps paying after the animal is gone. Removing an animal from $I$ only delays its arrival in $R$; it buys time rather than preventing contamination. Nor does targeting $R$ make the model indifferent to early removal: $\lambda_I$ inherits the full downstream cost of the cascade $I \to R \to W$ through the adjoint chain of Section~\ref{sec:adjoint-system}, so the optimizer values early removal without the control needing to be pointed at $I$.

One caveat is a limitation of the formulation rather than of the parameters. $R$ lumps pre-clinical shedders --- roughly 36 of its 53 weeks, and indistinguishable from healthy deer in the field --- with visibly clinical animals, so a control acting on all of $R$ presupposes test-and-cull. A purely sight-based cull would reach only the clinical fraction, about a third of $R$, and modelling it faithfully would require splitting $R$ in two.

The controlled state system is
\begin{align}
\dot S &= \Div\big(D_S\,\Grad S\big) + rS\Big(1-\frac{N}{K}\Big) - \beta_1 S R - \beta_2 S W, \label{eq:cS} \\
\dot I &= \Div\big(D_I\,\Grad I\big) + \beta_1 S R + \beta_2 S W - \sigma I, \label{eq:cI} \\
\dot R &= \Div\big(D_R\,\Grad R\big) + \sigma I - \alpha R \;-\; v\,R, \label{eq:cR} \\
\dot W &= \rho R - \delta W, \label{eq:cW}
\end{align}
on $\Gamma\times(0,T)$, with initial data
\begin{equation}
S(\cdot,0)=S_\infty,\quad I(\cdot,0)=I_0,\quad R(\cdot,0)=0,\quad W(\cdot,0)=0
\label{eq:ic}
\end{equation}
as in Section~\ref{sec:params-ic}, and the zero-flux conditions \eqref{eq:bc}, now written
\begin{equation}
\big(D_X\Grad X\big)\cdot\boldsymbol\nu = 0 \quad\text{on }\partial\Gamma\times(0,T), \qquad X\in\{S,I,R\}.
\label{eq:cbc}
\end{equation}
Note that \eqref{eq:ic} carries no dependence on $v$: the spin-up that produces $S_\infty$ runs with no infection present and hence with no culling, so the initial state is control-independent --- a fact used in Section~\ref{sec:discrete-adjoint}. Collecting the state as $y = (y_S,y_I,y_R,y_W) := (S,I,R,W)$, the reaction terms of \eqref{eq:cS}--\eqref{eq:cW}, \emph{including} the control term, are $f = (f_S,f_I,f_R,f_W)$ with
\begin{equation}
\begin{aligned}
f_S(y,v) &= G(S,I,R) - \beta_1 SR - \beta_2 SW, \qquad
&&G(S,I,R) := rS\Big(1-\frac{S+I+R}{K}\Big),\\
f_I(y,v) &= \beta_1 SR + \beta_2 SW - \sigma I, \qquad
&&f_R(y,v) = \sigma I - (\alpha+v) R,\\
f_W(y,v) &= \rho R - \delta W, &&
\end{aligned}
\label{eq:fdef}
\end{equation}
and, with $D_W := 0$, the state system reads compactly as
\begin{equation}
\dot y_k = \Div\big(D_k\Grad y_k\big) + f_k(y,v), \qquad k\in\{S,I,R,W\}.
\label{eq:state-compact}
\end{equation}

\subsection{Admissible controls and the objective functional}

Culling effort is bounded above by logistics and below by zero, so we restrict attention to
\begin{equation}
\Vad = \Big\{\,v\in L^\infty\big(\Gamma\times(0,T)\big)\;:\;0\le v(\mathbf{x},t)\le v_{\max}\ \text{ for a.e. }(\mathbf{x},t)\,\Big\},
\end{equation}
a closed, bounded, convex subset of $L^2(\Gamma\times(0,T))$, with $v_{\max}$ the largest sustained per-capita removal rate the management program can achieve.

\paragraph{Piecewise-constant-in-time controls.} The implementation optimizes over a closed convex \emph{subset} of $\Vad$: controls constant in time on each of $N_b$ blocks $[t_{b-1},t_b)$ partitioning $[0,T]$,
\begin{equation}
\Vad^{\,b} = \Big\{\,v\in\Vad \;:\; v(\mathbf{x},t) = v_b(\mathbf{x})\ \text{ for } t\in[t_{b-1},t_b),\ b=1,\dots,N_b \,\Big\} \subset \Vad,
\label{eq:Vadblock}
\end{equation}
with annual blocks ($N_b = 20$) by default. Because $\Vad^{\,b}\subset\Vad$, the attained minimum is an upper bound for the minimum over $\Vad$, and the restriction must be justified on its own terms. Two arguments do so. An unrestricted $v$ is free to change every $7.3$~days; no management agency re-tunes a cull at that cadence, so \eqref{eq:Vadblock} moves the admissible set \emph{toward} what can actually be enacted and the reported $J$ is a more honest estimate of an achievable cost. And the restriction is what makes the algorithm of Section~\ref{sec:algorithm} tractable: the discrete control has dimension $N_b\times N$, so at $N = 210{,}058$ nodes and $N_t = 1000$ steps an unrestricted control carries $2.1\times10^8$ degrees of freedom, and the limited-memory quasi-Newton method below --- which retains $2m+5$ vectors of that size --- would need some $39$~GiB at $m=10$. Annual blocking divides both by fifty. Shorter blocks are available at proportionate cost, and setting the block length to $\Delta t$ recovers $\Vad$ itself.

Nothing else in the derivation changes: since $v$ is constant across the steps of a block, the chain rule gives the block's reduced gradient as the \emph{sum} of the per-step gradients \eqref{eq:discrete-gradient} over that block, and weighting each block by its duration in the space-time inner product leaves $\nabla\hat J$ the Riesz representative it was before.

The objective functional to be minimized is
\begin{equation}
\boxed{\;J(v) \;=\; \int_0^T\!\!\int_\Gamma \Big[\,c_1\,R \;+\; c_2\,W \;+\; c_3\,v^2 \;+\; c_4\,v\,\Big]\,d\Gamma\,dt\;}
\label{eq:objective}
\end{equation}
with weights $c_1,c_2,c_4\ge 0$ and $c_3>0$, where $(S,I,R,W)$ solves \eqref{eq:cS}--\eqref{eq:cbc} driven by $v$. The four terms are, in order:
\begin{itemize}
\item $c_1 R$ --- the accumulated \emph{shedding} burden. Penalizing $R$ rather than $I$ matches both the control's target and what actually causes harm: an animal in $I$ transmits nothing and contaminates nothing. The upstream class is not thereby ignored, since a cost on $R$ propagates back to $I$ through the adjoint (Section~\ref{sec:adjoint-system}).
\item $c_2 W$ --- the accumulated environmental prion load. This term is what makes the CWD control problem qualitatively different from control of an acute directly-transmitted disease: contamination deposited early persists for years, so the optimizer is rewarded for acting \emph{early in the epizootic}, while there is still contamination to prevent, rather than merely for suppressing prevalence at any given instant.
\item $c_3 v^2$ --- the super-linear part of the intervention cost, whose marginal cost $2c_3v$ grows with the effort already being spent. The quadratic penalty makes $J$ strictly convex in $v$ for fixed state, keeps the optimal control a \emph{continuous} function of the adjoint rather than the bang-bang control a purely linear penalty would produce, and reflects the realistic diseconomy of scale in which doubling removal effort in one place costs more than twice as much.
\item $c_4 v$ --- the cost proportional to effort: so many dollars, or staff-days, per deer removed. This is the more literal description of what a cull costs, and, as \eqref{eq:pointwise-optimality} will show, it introduces a \emph{threshold}: the optimal effort is exactly zero wherever the local benefit $\lambda_R R$ falls below $c_4$, rather than the whisper of effort everywhere that a purely quadratic penalty produces. That threshold is the ``not here'' half of a cull-here-not-there map, and the half a management plan actually needs.
\end{itemize}
The combination is an \emph{elastic net} in $v$. Despite the $L^1$-like term, \eqref{eq:objective} is smooth on $\Vad$: because $v\ge0$ is imposed by the constraint set, $\int|v| = \int v$ there, so the kink of the absolute value sits exactly on the boundary of $\Vad$ and never in its interior. No proximal or splitting machinery is needed, and the projected methods of Section~\ref{sec:algorithm} apply unchanged. Setting $c_4=0$ recovers the purely quadratic objective, which is the default.

Since \eqref{eq:cS}--\eqref{eq:cbc} determines $y$ from $v$, we may regard $J$ as a function of $v$ alone. Writing $\hat J(v) := J(y(v),v)$ for the \emph{reduced} objective, the control problem is
\begin{equation}
\min_{v\in\Vad}\ \hat J(v).
\label{eq:ocp}
\end{equation}

\paragraph{Existence, and what is known about uniqueness.} A minimizer exists. $\Vad$ is bounded, closed and convex in $L^2(\Gamma\times(0,T))$; energy estimates on \eqref{eq:cS}--\eqref{eq:cbc} bound a minimizing sequence uniformly in $L^2(0,T;H^1(\Gamma))$ with time derivatives bounded in $L^2(0,T;H^1(\Gamma)^*)$, so an Aubin--Lions compactness argument upgrades weak convergence to strong convergence in $L^2$ --- which is what is needed to pass to the limit in the bilinear terms $SR$, $SW$, $SN/K$ and $vR$ --- while $c_3>0$ makes the control cost weakly lower semicontinuous. This is the argument of \citet[Thm.~A.2]{Neilan2011} for a parabolic SIR system of the same type.

Uniqueness is another matter. $\hat J$ is not convex in $v$: the state system is nonlinear in the terms just named, so $v\mapsto y(v)$ is nonlinear and \eqref{eq:ocp} may possess several local minimizers. \citet{Clayton2010}, for a raccoon-rabies SEIR model with a cost \emph{linear} in the control, establish existence together with the Pontryagin necessary conditions; the optimal control is bang-bang with a possible order-one singular arc, and no uniqueness is claimed. \citet{Neilan2011}, whose parabolic SIR model with a quadratic control cost on a box-constrained control is the closest published analogue of \eqref{eq:ocp}, do prove uniqueness of the solution of the \emph{optimality system}, and hence of the optimal control, but only for $T$ sufficiently small (their Thm.~A.5), through a Gronwall estimate whose admissible horizon is governed by the Lipschitz constants of the nonlinearities and is not quantified. That result is not available to us: our horizon is $T=20$~yr, chosen so that the environmental compartment's multi-year memory $1/\delta$ is resolved, which is the opposite of the small-time regime. The method of Section~\ref{sec:algorithm} accordingly locates a stationary point rather than a certified global minimizer, and re-running from several initial guesses remains good practice.

\subsection{Choosing the weights}

Only the ratios $c_1:c_2:c_3:c_4$ matter, since scaling $J$ by a positive constant leaves the minimizer unchanged. The weights carry units --- $c_1$ is (cost)$\cdot$km$^2$/(deer$\cdot$year), $c_2$ is (cost)$\cdot$km$^2$/(prion-unit$\cdot$year), $c_3$ is (cost)$\cdot$km$^2\cdot$year, $c_4$ is (cost)$\cdot$km$^2$ --- and should be chosen so that the integrals in \eqref{eq:objective} are of comparable magnitude at a representative control, otherwise the optimizer effectively ignores the smallest term. A practical calibration is to run the uncontrolled model once, record the integrals at $v\equiv \tfrac12 v_{\max}$, and rescale.

Applying that recipe gives the adopted $c_1 = 1$, $c_2 = 0.5$, $c_3 = 2$, all \tagGUESS{}. Fix $c_1 = 1$ as the numeraire. The environmental compartment settles at its quasi-equilibrium $W = \rho R/\delta = 2R$ exactly, so a unit weight on $W$ would give the environmental term precisely twice the leverage of the shedding term for the same underlying burden; $c_2 = 0.5$ puts them on equal footing. (That this ratio is exact rather than approximate is a small dividend of costing $R$ instead of $I$: $W$ and $R$ are related by a single linear ODE.) For $c_3$, at half effort the control term is $c_3 v^2 = 0.5$, against $c_1 R$ at the scale of the endemic shedding density $R^* = 0.70$ deer/km$^2$ --- the two within a factor of $1.4$, which is what ``comparable'' has to mean here. Without a weight of roughly this size the quadratic penalty is negligible against the state terms and the optimizer simply pins $v$ at $v_{\max}$ everywhere. Of the three, $c_3$ is the one most in need of replacement by a real per-deer cull cost set against a real valuation of a shedding deer.

The linear weight $c_4$ is set to $0$ by default, so the default configuration is the purely quadratic one. It is worth raising for two different reasons, calling for different values. If a genuine per-deer removal cost is available, $c_4$ is where it belongs, and its magnitude is fixed by the same numeraire as $c_1$. If instead the aim is a \emph{sparse} strategy --- a map that says where not to cull as well as where to --- then $c_4$ is a tuning parameter, and the threshold it sets in \eqref{eq:pointwise-optimality} should be read against the observed range of $\lambda_R R$. The motivation is concrete: on a five-year trial the purely quadratic objective returned a peak effort of $0.042\ \mathrm{yr}^{-1}$ against a space-time mean of $0.0045\ \mathrm{yr}^{-1}$, which reads as ``cull everywhere, barely'' and is not implementable as stated.

Raising $c_4$ while lowering $c_3$ has one real cost, and it is a conditioning cost rather than a modelling one. The quadratic weight contributes a uniform $2c_3\mathcal{I}$ floor to the reduced Hessian (Section~\ref{sec:lbfgs}); without it only the compact state-coupling operator remains, and its spectrum decays to zero. The limit $c_3\to0$ is the bang-bang problem, and is where convergence degrades --- a continuous dial rather than a cliff, but the reason $c_3>0$ is required rather than merely recommended. One tempting move should be refused: smoothing $|v|$ by a surrogate such as $\sqrt{v^2+\varepsilon}$ is unnecessary, since $v\ge0$ makes the term linear on $\Vad$, and actively harmful, since the surrogate's curvature is $\varepsilon^{-1/2}$ at the origin and decays like $\varepsilon/v^3$ away from it --- a spread of order $10^6$ across the box at $\varepsilon=10^{-6}$, manufacturing exactly the ill-conditioning it was meant to avoid.

The bound $v_{\max}$ deserves the same scrutiny, since it is what makes the box constraint bite. Being a per-capita removal rate on $R$, $v_{\max} = 1\ \mathrm{yr}^{-1}$ sustained for a year removes $1 - e^{-1} = 63\%$ of the shedding class --- at or beyond the ceiling of any real agency cull program, and beyond what the most intensive Wisconsin CWD-zone efforts sustained landscape-wide \citep{Wasserberg2009}. \tagGUESS{}.

% ======================================================================
\section{Derivation of the Adjoint System}
\label{sec:adjoint-derivation}
% ======================================================================

To apply any gradient-based method to \eqref{eq:ocp} we need $\hat J'(v)$. A na\"ive finite-difference gradient would require one state solve per control degree of freedom, which is hopeless at this dimension; the adjoint method obtains the entire gradient at the cost of one additional backward, \emph{linear} solve, regardless of how many control variables there are. The construction is the standard one for PDE-constrained optimization \citep{Hinze2009}, specialized to the surface reaction--diffusion system at hand.

\subsection{The Lagrangian and its derivative}

Introduce adjoint variables $\lambda = (\lambda_S,\lambda_I,\lambda_R,\lambda_W)$ and append the state equations to the objective:
\begin{equation}
\Lag(y,v,\lambda) \;=\; J(y,v) \;-\; \sum_{k}\int_0^T\!\!\int_\Gamma \lambda_k\Big[\dot y_k - \Div\big(D_k\Grad y_k\big) - f_k(y,v)\Big]\,d\Gamma\,dt,
\label{eq:lagrangian}
\end{equation}
the sum running over $k\in\{S,I,R,W\}$. For any $y$ satisfying the state equations the bracket vanishes and $\Lag = J$ for every $\lambda$; the point of the construction is that we may choose $\lambda$ to eliminate the unknown state sensitivities from the derivative. Name the running cost and the Hamiltonian density
\begin{equation}
\ell(y,v) := c_1 R + c_2 W + c_3 v^2 + c_4 v, \qquad
\Ham(y,v,\lambda) := \ell(y,v) + \sum_k \lambda_k f_k(y,v).
\label{eq:hamiltonian}
\end{equation}

Perturb the control, $v \to v + \varepsilon\,\delta v$ with $v+\varepsilon\,\delta v\in\Vad$, and let $\delta y$ be the induced state sensitivity. Linearizing \eqref{eq:state-compact},
\begin{equation}
\dot{\delta y_k} - \Div\big(D_k\Grad \delta y_k\big) - \sum_j \frac{\partial f_k}{\partial y_j}\,\delta y_j \;=\; \frac{\partial f_k}{\partial v}\,\delta v,
\label{eq:sensitivity}
\end{equation}
with $\delta y_k(\cdot,0)=0$, since the initial data \eqref{eq:ic} does not depend on the control, and the same zero-flux conditions \eqref{eq:cbc}; note $\partial f_k/\partial v = 0$ except for $\partial f_R/\partial v = -R$. Taking the G\^ateaux derivative of \eqref{eq:lagrangian} in the direction $\delta v$ and using \eqref{eq:sensitivity},
\begin{equation}
\delta\Lag = \int_0^T\!\!\int_\Gamma\Big[\frac{\partial \ell}{\partial v}\delta v + \sum_k\frac{\partial \ell}{\partial y_k}\delta y_k\Big]
\;-\;\sum_k\int_0^T\!\!\int_\Gamma\lambda_k\Big[\dot{\delta y_k} - \Div\big(D_k\Grad\delta y_k\big) - \sum_j\frac{\partial f_k}{\partial y_j}\delta y_j - \frac{\partial f_k}{\partial v}\delta v\Big].
\label{eq:dLag}
\end{equation}
Two integrations by parts move all derivatives off $\delta y$. In time,
\[
\int_0^T \lambda_k\,\dot{\delta y_k}\,dt \;=\; \lambda_k(T)\,\delta y_k(T) - \int_0^T \dot\lambda_k\,\delta y_k\,dt,
\]
the $t=0$ term vanishing because $\delta y_k(0)=0$, and the remaining boundary term removed by imposing the \emph{terminal} condition
\begin{equation}
\lambda_k(\cdot,T) = 0, \qquad k\in\{S,I,R,W\}.
\label{eq:terminal}
\end{equation}
This is the origin of the backward-in-time character of the adjoint system: the state carries an initial condition, so the co-state must carry a terminal one. In space, applying the surface divergence theorem twice with the zero-flux conditions \eqref{eq:cbc} on $\delta y_k$ and the same conditions imposed on $\lambda_k$,
\[
-\int_\Gamma \lambda_k\,\Div\big(D_k\Grad\delta y_k\big)\,d\Gamma
= \int_\Gamma \Grad \delta y_k\cdot D_k^{\mathsf T}\Grad\lambda_k \,d\Gamma
= -\int_\Gamma \delta y_k\,\Div\big(D_k^{\mathsf T}\Grad\lambda_k\big)\,d\Gamma.
\]
Because $D_k^{\mathsf T}=D_k$ (Section~\ref{sec:diffusion}), the adjoint equations carry the \emph{same} diffusion operator as the state equations, not its transpose --- which is what will allow the discrete forward and backward sweeps to share a single matrix factorization. Substituting both integrations by parts into \eqref{eq:dLag} and grouping the terms multiplying each $\delta y_k$,
\begin{equation}
\delta\Lag = \sum_k\int_0^T\!\!\int_\Gamma \delta y_k\,
\underbrace{\Big[\dot\lambda_k + \Div\big(D_k\Grad\lambda_k\big) + \frac{\partial \Ham}{\partial y_k}\Big]}_{\text{choose }\lambda\text{ to kill this}}
\;+\;\int_0^T\!\!\int_\Gamma \frac{\partial\Ham}{\partial v}\,\delta v,
\label{eq:collected}
\end{equation}
where $\partial\Ham/\partial y_k = \partial \ell/\partial y_k + \sum_j \lambda_j\,\partial f_j/\partial y_k$ and $\partial\Ham/\partial v = 2c_3 v + c_4 - \lambda_R R$. Choosing $\lambda$ so that the underbraced bracket vanishes leaves a formula for the derivative of the reduced objective involving \emph{only} $\delta v$.

\subsection{The adjoint system}
\label{sec:adjoint-system}

Setting the bracket in \eqref{eq:collected} to zero gives
\begin{equation}
-\dot\lambda_k - \Div\big(D_k\Grad\lambda_k\big) \;=\; \frac{\partial \ell}{\partial y_k} + \sum_{j}\lambda_j\,\frac{\partial f_j}{\partial y_k},
\qquad \lambda_k(\cdot,T)=0.
\label{eq:adjoint-general}
\end{equation}
Abbreviating the partial derivatives of the logistic growth term of \eqref{eq:fdef} as
\begin{equation}
G_S := r\Big(1-\frac{N}{K}\Big) - \frac{rS}{K},
\qquad
G_I = G_R := -\frac{rS}{K},
\label{eq:Gderivs}
\end{equation}
the Jacobian of $f$ is
\begin{equation}
\left[\frac{\partial f_j}{\partial y_k}\right]_{j,k}
=
\begin{pmatrix}
G_S - \beta_1 R - \beta_2 W & G_I & G_R - \beta_1 S & -\beta_2 S\\[2pt]
\beta_1 R + \beta_2 W & -\sigma & \beta_1 S & \beta_2 S\\[2pt]
0 & \sigma & -\alpha - v & 0\\[2pt]
0 & 0 & \rho & -\delta
\end{pmatrix},
\label{eq:jacobian}
\end{equation}
with rows indexed by $j$ (which equation) and columns by $k$ (with respect to which state). The sum in \eqref{eq:adjoint-general} is precisely the $k$-th entry of $\big[\partial f/\partial y\big]^{\mathsf T}\lambda$: \emph{the adjoint system is driven by the transpose of the state Jacobian}, which is why its couplings run in the opposite direction. Note also where the control lands in \eqref{eq:jacobian} --- on the diagonal alone, the $(R,R)$ entry gaining $-v$, with every off-diagonal coupling untouched. Finally $\partial\ell/\partial R = c_1$, $\partial\ell/\partial W = c_2$, and $\partial\ell/\partial S = \partial\ell/\partial I = 0$. Substituting, the adjoint system is
\begin{align}
-\dot\lambda_S - \Div\big(D_S\Grad\lambda_S\big) &= \lambda_S\big[G_S - \beta_1 R - \beta_2 W\big] + \lambda_I\big[\beta_1 R + \beta_2 W\big], \label{eq:adjS}\\
-\dot\lambda_I - \Div\big(D_I\Grad\lambda_I\big) &= \lambda_S\,G_I - \lambda_I\,\sigma + \lambda_R\,\sigma, \label{eq:adjI}\\
-\dot\lambda_R - \Div\big(D_R\Grad\lambda_R\big) &= c_1 + \lambda_S\big[G_R - \beta_1 S\big] + \lambda_I\,\beta_1 S - \lambda_R\big(\alpha + v\big) + \lambda_W\,\rho, \label{eq:adjR}\\
-\dot\lambda_W &= c_2 - \lambda_S\,\beta_2 S + \lambda_I\,\beta_2 S - \lambda_W\,\delta, \label{eq:adjW}
\end{align}
on $\Gamma\times(0,T)$, with $\lambda_k(T)=0$ and $\big(D_X\Grad\lambda_X\big)\cdot\boldsymbol\nu=0$ on $\partial\Gamma$ for $X\in\{S,I,R\}$.

Several features are worth drawing out.
\begin{itemize}
\item \textbf{It is linear, and it runs backward.} Although the state system is nonlinear, \eqref{eq:adjS}--\eqref{eq:adjW} is linear in $\lambda$, the state and control entering only as known coefficients, so there is no Newton iteration in the backward solve. Substituting $\tau = T-t$ turns it into a forward-in-time linear reaction--diffusion system of exactly the same structural type as the state system, so the IMEX machinery of Section~\ref{sec:fem} applies verbatim --- and $\lambda_W$, like $W$, carries no diffusion and is advanced nodewise.

\item \textbf{The couplings are transposed.} Where the state propagates infection forward along $S\to I\to R\to W$, the adjoint propagates \emph{cost} backward along the same chain: $\lambda_W$ feeds $\lambda_R$ through $\lambda_W\rho$ in \eqref{eq:adjR}, $\lambda_R$ feeds $\lambda_I$ through $\lambda_R\sigma$ in \eqref{eq:adjI}, and so on.

\item \textbf{The cost weights are sources.} The running cost enters only as the constant forcings $c_1$ in \eqref{eq:adjR} and $c_2$ in \eqref{eq:adjW}; with $c_1=c_2=0$ we would obtain $\lambda\equiv 0$ and a zero gradient. Both sources sit in the \emph{downstream} half of the chain, so nothing forces $\lambda_S$ or $\lambda_I$ directly and they become nonzero only by inheritance.

\item \textbf{Interpretation.} $\lambda_k(\mathbf{x},t)$ is the marginal cost, over the remaining horizon $[t,T]$, of one additional unit of compartment $k$ --- a shadow price. In particular $\lambda_R$ measures the total future damage caused by one more shedding deer: its own accrued cost $c_1$, the infections it will go on to cause directly, and the environmental contamination it will deposit and that will then persist for a mean $1/\delta$ years. That $\lambda_R$ is exactly the multiplier appearing in the gradient below is the mathematical content of the intuition that one should cull hardest where and when a shedding animal is most costly.

\item \textbf{Early removal is valued without being costed.} Equation~\eqref{eq:adjI} carries the term $+\lambda_R\sigma$, so $\lambda_I$ is driven by $\lambda_R$ through exactly the progression rate at which $I$ becomes $R$; since $\sigma$ is also the only exit from $I$, a slowly varying adjoint gives $\lambda_I \approx \lambda_R$. A sub-clinical animal is thus valued at the full shadow price of the shedding animal it is certain to become. Were a background mortality reinstated, this would become $\lambda_I \approx \lambda_R\,\sigma/(\sigma+\mu)$ --- the same price discounted by the probability of living to shed.

\item \textbf{Terminal payoff.} If a terminal cost $\int_\Gamma \Phi(y(\cdot,T))\,d\Gamma$ is added to \eqref{eq:objective} --- for example $c_T\int_\Gamma W(T)$ to penalize the contaminated landscape left behind at the horizon, which for CWD is arguably the burden that most deserves a terminal penalty --- then only the terminal condition changes, to $\lambda_k(\cdot,T) = \partial\Phi/\partial y_k$.

\item \textbf{The water branch.} The implementation replaces $G$ by $-rS$ wherever $K(\mathbf{x})<0.01$ (Section~\ref{sec:monolithic}). The adjoint must inherit exactly the same branch: there $G_S = -r$ and $G_I=G_R=0$. Failing to differentiate the conditional consistently is a common and silent source of wrong gradients.
\end{itemize}

The weak form used for the finite element solve is obtained exactly as in Section~\ref{sec:weak}, by multiplying \eqref{eq:adjS}--\eqref{eq:adjR} by $\psi\in V_h$ and applying the surface divergence theorem, with \eqref{eq:adjW} advanced nodewise. Its bilinear (diffusion plus mass) parts are \emph{identical} to those of the state problem, term for term.

% ======================================================================
\section{Optimality Condition and the Reduced Gradient}
\label{sec:optimality}
% ======================================================================

With $\lambda$ solving \eqref{eq:adjS}--\eqref{eq:adjW}, the bracket in \eqref{eq:collected} vanishes and we are left with
\begin{equation}
\boxed{\;
\hat J'(v)\,[\delta v] \;=\; \int_0^T\!\!\int_\Gamma \Big(2c_3\,v + c_4 - \lambda_R\,R\Big)\,\delta v \;d\Gamma\,dt,
\qquad\text{i.e.}\qquad
\nabla \hat J(v) \;=\; 2c_3\,v + c_4 - \lambda_R\,R \;}
\label{eq:gradient}
\end{equation}
the gradient understood as the Riesz representative in $L^2(\Gamma\times(0,T))$. This is the payoff of the whole construction: \emph{one} state solve plus \emph{one} adjoint solve yield the gradient with respect to every control degree of freedom simultaneously.

Equation~\eqref{eq:gradient} has a clean reading. The optimizer increases culling wherever the product $\lambda_R R$ --- the density of shedding animals times the marginal future cost of a shedding animal --- exceeds the marginal cost $2c_3 v + c_4$ of effort at the level already being spent there. The product form matters: effort is worthless where there is nothing to cull ($R \to 0$) \emph{and} where culling buys nothing ($\lambda_R \to 0$, as it does near $t=T$, when there is no remaining horizon over which the prevented contamination could do damage). The two cost weights enter differently: the quadratic term contributes $2c_3v$, which vanishes as the effort does and so never by itself argues for spending nothing, while the linear term contributes the constant $c_4$, which does not, and so sets a floor on how valuable a location must be before any effort is worth spending there.

\subsection{First-order necessary condition}

Since $\Vad$ is convex, a local minimizer $v^*$ satisfies the variational inequality
\begin{equation}
\int_0^T\!\!\int_\Gamma \Big(2c_3 v^* + c_4 - \lambda_R R\Big)\big(w - v^*\big)\,d\Gamma\,dt \;\ge\; 0
\qquad \text{for all } w\in\Vad,
\label{eq:vi}
\end{equation}
with $(S,I,R,W)$ and $\lambda$ the state and adjoint associated with $v^*$. Because the constraint acts pointwise, \eqref{eq:vi} is equivalent to
\begin{equation}
v^*(\mathbf{x},t)
\;=\; P_{\Vad}\!\left(\frac{\lambda_R R - c_4}{2c_3}\right)
\;=\; \min\left\{\,v_{\max},\ \max\left\{\,0,\ \frac{\lambda_R R - c_4}{2c_3}\,\right\}\right\},
\label{eq:pointwise-optimality}
\end{equation}
$P_{\Vad}$ denoting pointwise projection onto $[0,v_{\max}]$. This is where the two control weights show their different characters. The lower bound is active --- the optimum is \emph{exactly} zero, not merely small --- on the whole set where $\lambda_R R \le c_4$, so $c_4$ acts as a threshold on the local benefit and produces genuinely sparse strategies; above that threshold the effort ramps linearly in $\lambda_R R$ at slope $1/(2c_3)$. At $c_4=0$ the threshold collapses to the single point $\lambda_R R = 0$ and the familiar $v^* = P_{\Vad}(\lambda_R R/2c_3)$ is recovered. Equivalently,
\begin{equation}
v^* = P_{\Vad}\big(v^* - \alpha\,\nabla\hat J(v^*)\big) \quad\text{for any } \alpha>0,
\label{eq:stationarity}
\end{equation}
which is the stationarity measure the code monitors (with $\alpha=1$) for its convergence test.

Formula~\eqref{eq:pointwise-optimality} suggests the classical \emph{forward--backward sweep} of \citet{Lenhart2007}: solve the state forward, the adjoint backward, set $v$ to the right-hand side, repeat. That iteration is simple but has no descent guarantee and is known to diverge when the coupling is strong --- which for this model it is, since $\lambda_R$ depends on $v$ through the whole $-\lambda_R(\alpha+v)$ term. The method of the next section keeps the same forward--backward structure but replaces the naive update with a safeguarded descent step, at the cost of a few extra forward solves per iteration.

% ======================================================================
\section{Numerical Solution of the Control Problem}
\label{sec:control-numerics}
% ======================================================================

\subsection{The discrete adjoint of the IMEX scheme}
\label{sec:discrete-adjoint}

There are two ways to obtain a discrete gradient: discretize the continuous adjoint (\emph{optimize-then-discretize}), or differentiate the discrete state system exactly (\emph{discretize-then-optimize}). We take the latter. The two agree to the order of the discretization, but only the discrete adjoint yields the \emph{exact} gradient of the quantity the computer actually minimizes. That exactness matters in practice: an approximate gradient is not a descent direction near a stationary point, and the Armijo line search of Section~\ref{sec:algorithm} then stalls for reasons that look like a bug in the model rather than an inconsistency in the derivative.

Let $\mathbf{y}^n\in\mathbb{R}^{4N}$ be the nodal coefficient vector of all four compartments at $t^n=n\Delta t$, and $\mathbf{v}^n\in\mathbb{R}^{N}$ the nodal control on the slab $[t^n,t^{n+1})$, with $N_t = \lfloor T/\Delta t\rfloor$. The IMEX step \eqref{eq:imex}, assembled monolithically as in Section~\ref{sec:monolithic}, reads
\begin{equation}
B\,\mathbf{y}^{n+1} \;=\; \mathbf{L}\big(\mathbf{y}^n,\mathbf{v}^n\big), \qquad n=0,\dots,N_t-1,
\label{eq:discrete-state}
\end{equation}
with $\mathbf{L}$ the nonlinear load vector of Section~\ref{sec:monolithic} including the control term $-\Delta t\,v^n R^n$ in the $R$-row. The discrete objective, using a right-endpoint rule for the state terms and a left-endpoint rule for the control terms, is
\begin{equation}
J_h(\mathbf{v}) \;=\; \Delta t\sum_{n=1}^{N_t}\int_{\Gamma_h}\big(c_1 R_h^n + c_2 W_h^n\big)\,dA
\;+\; \Delta t\sum_{n=0}^{N_t-1}\int_{\Gamma_h} \big(c_3\,(v_h^n)^2 + c_4\,v_h^n\big)\,dA .
\label{eq:discrete-objective}
\end{equation}
Introducing one multiplier vector $\boldsymbol\lambda^{n+1}\in\mathbb{R}^{4N}$ per time step,
\begin{equation}
\Lag_h \;=\; J_h(\mathbf{v}) \;-\; \sum_{n=0}^{N_t-1}\big(\boldsymbol\lambda^{n+1}\big)^{\mathsf T}\Big(B\,\mathbf{y}^{n+1} - \mathbf{L}(\mathbf{y}^n,\mathbf{v}^n)\Big),
\end{equation}
and letting $\mathbf{g}\in\mathbb{R}^{4N}$ be assembled from $\int_{\Gamma_h}(c_1\psi_R + c_2\psi_W)\,dA$, so that $\partial J_h/\partial \mathbf{y}^n = \Delta t\,\mathbf{g}$ for $1\le n\le N_t$, stationarity of $\Lag_h$ with respect to $\mathbf{y}^{N_t}$ and to $\mathbf{y}^{n}$, $1\le n\le N_t-1$, gives
\begin{equation}
B^{\mathsf T}\boldsymbol\lambda^{N_t} = \Delta t\,\mathbf{g},
\qquad
B^{\mathsf T}\boldsymbol\lambda^{n} \;=\; \left[\frac{\partial \mathbf{L}}{\partial\mathbf{y}}\big(\mathbf{y}^n,\mathbf{v}^n\big)\right]^{\mathsf T}\boldsymbol\lambda^{n+1} \;+\; \Delta t\,\mathbf{g}.
\label{eq:discrete-adjoint}
\end{equation}
No equation is needed for $\mathbf{y}^0$, which is fixed by the initial condition; because that condition is control-independent (Section~\ref{sec:spinup}), $\partial\mathbf{y}^0/\partial\mathbf{v}=0$ and no corresponding term enters the reduced gradient.

Since $B^{\mathsf T}=B$ (Section~\ref{sec:fem}), \eqref{eq:discrete-adjoint} may be solved with the \emph{same} block factorizations \eqref{eq:blockdiag} already computed for the forward sweep: the backward solve costs one vector assembly and one back-substitution per step, exactly like the forward solve, and no second matrix is assembled or factorized anywhere in the code. The transposed Jacobian $[\partial\mathbf{L}/\partial\mathbf{y}]^{\mathsf T}\boldsymbol\lambda^{n+1}$ is assembled row-wise by differentiating the integrand of Section~\ref{sec:monolithic} with respect to each state, using $G_S$, $G_I=G_R$ of \eqref{eq:Gderivs} and inheriting the $K<0.01$ branch; the resulting terms are exactly the right-hand sides of \eqref{eq:adjS}--\eqref{eq:adjW} minus their $c_1,c_2$ forcings, which \eqref{eq:discrete-adjoint} supplies separately through $\Delta t\,\mathbf{g}$. The discrete adjoint is thus recognizably a discretization of the continuous adjoint --- but it is derived, not guessed, and it is exact.

Stationarity of $\Lag_h$ with respect to $\mathbf{v}^n$ gives the discrete gradient
\begin{equation}
\frac{\partial J_h}{\partial \mathbf{v}^n} \;=\; \Delta t\Big[\,2c_3\,M\,\mathbf{v}^n \;+\; c_4\,M\mathbf{1} \;-\; \mathbf{m}\big(\mathbf{y}^n,\boldsymbol\lambda^{n+1}\big)\Big],
\qquad
\mathbf{m}_j = \int_{\Gamma_h} R_h^n\,\lambda_{R,h}^{n+1}\,\varphi_j\,dA ,
\label{eq:discrete-gradient}
\end{equation}
with $\mathbf{1}$ the vector of ones. Dividing by $M_L$ converts this from an element of the dual into the $L^2$ Riesz representative used by the line search,
\begin{equation}
g^n \;:=\; \frac{1}{\Delta t}M_L^{-1}\frac{\partial J_h}{\partial\mathbf{v}^n}
\;\approx\; 2c_3\,v^n + c_4 - R^n\,\lambda_R^{n+1},
\end{equation}
the discrete counterpart of \eqref{eq:gradient}. The corresponding inner product on space--time control fields,
\begin{equation}
\big\langle a, b\big\rangle \;=\; \Delta t\sum_{n=0}^{N_t-1}\sum_{i=1}^{N} \big(M_L\big)_{ii}\,a^n_i\,b^n_i,
\qquad \big(M_L\big)_{ii} = \int_{\Gamma_h}\varphi_i\,dA,
\label{eq:inner-product}
\end{equation}
is a quadrature for $\int_0^T\int_\Gamma ab$ and is the pairing in which \eqref{eq:discrete-gradient} is exactly the derivative of \eqref{eq:discrete-objective}.

\paragraph{Matching quadrature rules.} For \eqref{eq:discrete-adjoint} to be the exact transpose, every term of the adjoint load vector must be assembled with the \emph{same} quadrature rule as the term of $\mathbf{L}$ it differentiates. Quadrature is a linear functional, so differentiating a quadrature sum and quadrature-summing the analytic derivative agree exactly --- but only when the rule is the same; if each form is allowed to pick its own automatically estimated degree, the assembled adjoint is only approximately transposed and the Taylor test of Section~\ref{sec:verification} degrades from second order to first. Two rules appear, paired consistently: the mass terms of $\mathbf{L}$ and their transpose both use the vertex rule that produces $M_L$, while all reaction, cost and gradient terms use one explicitly pinned Gauss degree, shared across the state load vector, the adjoint load vector, the running cost and the gradient form. The quadrature of \eqref{eq:bilinear} is unconstrained by this requirement, since $B$ appears identically on the left of the forward and backward steps and never enters the transpose identity.

\subsection{Projected gradient descent with an Armijo line search}
\label{sec:algorithm}

Two optimizers consume the reduced gradient, selected at run time. This subsection describes the projected-gradient method, the simpler of the two and the reference implementation; Section~\ref{sec:lbfgs} describes the quasi-Newton method used by default. Both are fed the identical gradient.

The bound constraints are handled by projection. Writing $P$ for the pointwise projection onto $[0,v_{\max}]$ and $g^{(k)} = \nabla\hat J(v^{(k)})$, the update is $v^{(k+1)} = P(v^{(k)} - \alpha_k\,g^{(k)})$ with $\alpha_k$ chosen by backtracking. The sufficient-decrease (Armijo) test is imposed in the form appropriate to a projected method: with $d(\alpha) := P(v^{(k)}-\alpha g^{(k)}) - v^{(k)}$, accept the first $\alpha$ in the sequence $\alpha,\tau\alpha,\tau^2\alpha,\dots$ for which
\begin{equation}
\boxed{\;\hat J\big(v^{(k)} + d(\alpha)\big) \;\le\; \hat J\big(v^{(k)}\big) \;+\; \eta\,\big\langle g^{(k)},\,d(\alpha)\big\rangle\;}
\label{eq:armijo}
\end{equation}
with $\eta\in(0,1)$ (typically $10^{-4}$) and $\tau\in(0,1)$ (typically $\tfrac12$). Two facts make this the right test. Because $P$ projects onto a convex set, $\langle g, P(v-\alpha g)-v\rangle \le 0$ always, with equality if and only if $v$ satisfies \eqref{eq:stationarity}, so the right-hand side of \eqref{eq:armijo} is a genuine decrease target whenever $v^{(k)}$ is not already stationary. And when no bound is active the projection is the identity, $d(\alpha) = -\alpha g^{(k)}$, and \eqref{eq:armijo} collapses to the classical Armijo condition. Standard theory \citep{Nocedal2006} then gives that the backtracking terminates in finitely many steps and that every accumulation point of $\{v^{(k)}\}$ is stationary. Each evaluation of $\hat J$ costs one forward solve; the complete method is Algorithm~\ref{alg:sweep}.

\begin{figure}[h]
\centering
\fbox{\begin{minipage}{0.92\textwidth}
\textbf{Algorithm 1: Forward--backward sweep with projected Armijo descent.}
\smallskip
\hrule
\smallskip
Given $v^{(0)}\in\Vad$, initial step $\alpha_0>0$, parameters $\eta\in(0,1)$, $\tau\in(0,1)$, growth factor $\gamma\ge1$, tolerances $\epsilon_g,\epsilon_J$.
\begin{enumerate}
\item \textbf{Initial forward solve.} Advance \eqref{eq:discrete-state} from $\mathbf{y}^0$ with $v^{(0)}$, storing the entire trajectory $\{\mathbf{y}^n\}_{n=0}^{N_t}$, and evaluate $\hat J(v^{(0)})$ via \eqref{eq:discrete-objective}.
\item \textbf{For} $k=0,1,2,\dots$ until convergence:
\begin{enumerate}
\item \textbf{Backward solve.} Set $\boldsymbol\lambda^{N_t}$ from \eqref{eq:discrete-adjoint}. For $n = N_t-1,\dots,1$, evaluate the gradient slab $g^{(k),n} = 2c_3 v^{(k),n} + c_4 - R^n\lambda_R^{n+1}$ and then solve \eqref{eq:discrete-adjoint} for $\boldsymbol\lambda^{n}$, reusing the factors of $B$. Evaluate the last slab $g^{(k),0}$ from $\boldsymbol\lambda^{1}$.
\item \textbf{Stationarity test.} If $\lVert P(v^{(k)} - g^{(k)}) - v^{(k)}\rVert \le \epsilon_g$, \textbf{stop}.
\item \textbf{Line search.} Set $\alpha \leftarrow \gamma\alpha_k$. Repeat: form $v_{\text{trial}} = P(v^{(k)} - \alpha g^{(k)})$, run a forward solve to get $\hat J(v_{\text{trial}})$, and accept if \eqref{eq:armijo} holds; otherwise set $\alpha\leftarrow\tau\alpha$ and repeat. If $\alpha$ underflows, \textbf{stop} (line search failure).
\item \textbf{Update.} $v^{(k+1)} \leftarrow v_{\text{trial}}$, $\alpha_{k+1}\leftarrow\alpha$. If the relative decrease in $\hat J$ is below $\epsilon_J$, \textbf{stop}.
\end{enumerate}
\end{enumerate}
\end{minipage}}
\caption{The forward--backward sweep. Each outer iteration costs one backward solve plus $1+(\text{number of backtracks})$ forward solves, all reusing the block factorizations of $B$ computed once at start-up.}
\label{alg:sweep}
\end{figure}

An outer iteration costs one backward sweep and, typically, one to three forward sweeps. The dominant memory cost is the stored forward trajectory, since the backward sweep evaluates the Jacobian at $\mathbf{y}^n$: that is $4N(N_t+1)$ doubles, roughly $1.6\,$GB for a $10^5$-node mesh and $500$ steps. Where that is prohibitive the trajectory can be memory-mapped to disk; the classical alternative is checkpointing, storing every $m$-th state and recomputing the intermediate ones during the backward pass, which trades a factor $m$ in memory for roughly one extra forward solve. Finally, the magnitude of $\nabla\hat J$ is unit-dependent and not known in advance, so rather than fix $\alpha_0$ we size the first trial step so that $-\alpha_0 g^{(0)}$ moves the most gradient-sensitive node across the whole admissible range, $\alpha_0 = (v_{\max}-v_{\min})/\lVert g^{(0)}\rVert_\infty$; thereafter $\alpha$ is grown by $\gamma$ at the start of each line search and shrunk by $\tau$ within it.

\subsection{The default optimizer: L-BFGS-B}
\label{sec:lbfgs}

Projected gradient descent converges \emph{linearly}, at a rate governed by the conditioning of the reduced Hessian $\nabla^2\hat J$. That Hessian is $2c_3\mathcal{I}$ plus a compact operator arising through the state equation, so its spectrum is a cluster near $2c_3$ together with a tail of larger eigenvalues, and the method makes rapid initial progress and then crawls. Since each iteration costs a full forward--backward pair over $N_t$ steps, the iteration count is the whole story.

The bound constraint in \eqref{eq:Vadblock} is a \emph{simple box}, precisely the structure L-BFGS-B is designed for. It maintains an implicit approximation to $\nabla^2\hat J$ built from the last $m$ pairs $(s^{(k)},z^{(k)}) = (v^{(k+1)}-v^{(k)},\,g^{(k+1)}-g^{(k)})$, identifies the active bounds by a gradient-projection search, and takes a subspace quasi-Newton step on the remaining free variables; on the inactive set convergence is superlinear rather than linear, at unchanged cost per iteration. The implementation delegates to SciPy's L-BFGS-B, passing the reduced objective and gradient of Section~\ref{sec:discrete-adjoint} as a single value-and-gradient callable and the box $[0,v_{\max}]$ as bounds. Nothing in the discretization, the adjoint, or the gradient changes.

Three practical points follow. L-BFGS-B stores $2m+5$ vectors the size of the control, which is exactly why the blocking \eqref{eq:Vadblock} is not optional --- with an unrestricted control on the reference mesh those vectors total some $39$~GiB at the default $m=10$, against roughly $0.8$~GiB under annual blocking --- so $m$ trades directly against the block length and should be reduced before the blocks are shortened. SciPy's \texttt{ftol} is the same relative-decrease test the projected-gradient loop applies, but its \texttt{gtol} is the max-norm of the projected gradient, whereas the stationarity measure reported in Algorithm~\ref{alg:sweep} is the mass-weighted $L^2$ quantity $\lVert P(v-g)-v\rVert$; the two differ in scale and should not be compared directly. And L-BFGS-B may return a point other than the last one evaluated, so the driver performs one final forward solve at the returned control before writing any output rather than reusing whatever trajectory happened to be in memory.

Neither method escapes the non-convexity of $\hat J$ noted after \eqref{eq:ocp}: both locate a stationary point, not a certified global minimizer.

\subsection{Verification: the Taylor test}
\label{sec:verification}

Adjoint implementations fail silently: a gradient wrong by a sign, a missing Jacobian term, or an undifferentiated conditional still produces plausible-looking output, merely converging slowly or to the wrong control. The standard check is a Taylor test: for a fixed direction $d$ and $\varepsilon\to0$,
\begin{equation}
R_0(\varepsilon) := \big|\hat J(v+\varepsilon d) - \hat J(v)\big| = O(\varepsilon),
\qquad
R_1(\varepsilon) := \big|\hat J(v+\varepsilon d) - \hat J(v) - \varepsilon\,\langle \nabla\hat J(v), d\rangle\big| = O(\varepsilon^2).
\end{equation}
The second-order decay of $R_1$ holds only if the gradient is correct; any error leaves a term linear in $\varepsilon$ and degrades $R_1$ to first order. The test must be run at a control strictly inside the bounds, so that no projection is active and $\hat J$ is genuinely differentiable along $d$, and only over the range of $\varepsilon$ where round-off has not yet taken over.

% ======================================================================
\section{Implementation of the Control Solver}
\label{sec:control-implementation}
% ======================================================================

The control solver shares the state solver's mesh-generation workflow, parameter file and land-cover mappings; Table~\ref{tab:control-notation} extends Table~\ref{tab:code-notation} with the control-specific quantities. Three modules are involved.

\begin{itemize}
\item \texttt{cwd\_control\_problem.py} defines \texttt{CWDControlProblem}, which reads the mesh bundle, builds the mixed space and the land-cover fields, assembles the diffusion tensor \eqref{eq:Dtens} and the operator $B$, factorizes $B$ once, and compiles all of the variational forms of Table~\ref{tab:control-notation}. Its two workhorse methods are \texttt{forward(control, trajectory)}, which advances \eqref{eq:discrete-state} while accumulating \eqref{eq:discrete-objective} and recording the state history, and \texttt{backward(control, trajectory)}, which performs the sweep \eqref{eq:discrete-adjoint} and returns the gradient array. Placing both in one class is what guarantees that the forward and adjoint forms see the same mesh, the same coefficients, and --- critically --- the same quadrature rule.

\item \texttt{CWD\_optimal\_control.py} is the driver: it parses the mesh bundle flags identically to \texttt{CWD\_solver.py}, runs Algorithm~\ref{alg:sweep}, reports the cost history, and writes the optimal state, the optimal control and optionally the adjoint fields as XDMF/HDF5 time series. It also implements the Taylor test of Section~\ref{sec:verification} under \texttt{-{}-gradient-check}.

\item \texttt{utilities/susceptible\_spinup.py} implements the disease-free relaxation \eqref{eq:spinup-step}--\eqref{eq:spinup-stop} together with its on-disk cache and the signature check that invalidates it. It is called by \texttt{CWDControlProblem} and \texttt{CWD\_solver.py} alike, so both drivers start from the identical susceptible field; \texttt{-{}-recompute-equilibrium} forces the relaxation to be redone and \texttt{-{}-susceptible-equilibrium} points at a cache elsewhere than the mesh folder.
\end{itemize}

The control is stored as an array of shape $(N_b, N)$ --- one CG1 nodal field per block of \eqref{eq:Vadblock}, piecewise constant in time --- and saved as \texttt{optimal\_control.npy}, from which a subsequent run can be warm-started with \texttt{-{}-initial-control}. The per-iteration cost, gradient norm, step size and backtrack count are written to \texttt{optimization\_history.json}.

\begin{table}[h]
\centering
\small
\caption{Correspondence between the symbols of Sections~\ref{sec:control}--\ref{sec:control-numerics} and the implementation.}
\label{tab:control-notation}
\begin{tabular}{@{}lll@{}}
\toprule
Math & Code identifier & Role \\
\midrule
$v(\mathbf{x},t)$ & \texttt{control}, \texttt{v\_func} & culling control (array of blocks; CG1 field for one block) \\
$\lambda_S,\lambda_I,\lambda_R,\lambda_W$ & \texttt{Lam} (mixed \texttt{Function}) & adjoint state $\boldsymbol\lambda^{n+1}$ \\
$c_1$ & \texttt{cost\_shedding} & weight on the shedding burden $R$ \\
$c_2$ & \texttt{cost\_environment} & weight on the environmental prion load \\
$c_3$ & \texttt{cost\_control} & quadratic weight on the culling effort \\
$c_4$ & \texttt{cost\_control\_l1} & linear control weight; default $0$ \\
$v_{\min},\ v_{\max}$ & \texttt{control\_minimum}, \texttt{control\_maximum} & bounds of $\Vad$ \\
$N_b$ & \texttt{control\_block\_years} & control block length in \eqref{eq:Vadblock} \\
$m$ & \texttt{lbfgs\_memory} & L-BFGS-B history length, Section~\ref{sec:lbfgs} \\
$B$ & \texttt{A} & implicit operator \eqref{eq:bilinear}, factorized once \\
$\mathbf{L}(\mathbf{y}^n,\mathbf{v}^n)$ & \texttt{state\_rhs\_form} & forward load vector \eqref{eq:discrete-state} \\
$[\partial\mathbf{L}/\partial\mathbf{y}]^{\mathsf T}\boldsymbol\lambda + \Delta t\mathbf{g}$ & \texttt{adjoint\_rhs\_form} & adjoint load vector \eqref{eq:discrete-adjoint} \\
$\Delta t\,\mathbf{g}$ & \texttt{terminal\_rhs\_form} & terminal adjoint \eqref{eq:discrete-adjoint} \\
$\partial J_h/\partial\mathbf{v}^n$ & \texttt{gradient\_form} & reduced gradient \eqref{eq:discrete-gradient} \\
$M_L$ & \texttt{lumped\_mass} & lumped mass, Riesz map and inner product \eqref{eq:inner-product} \\
$\eta,\ \tau,\ \gamma$ & \texttt{armijo\_sufficient\_decrease}, & line search parameters \\
 & \texttt{armijo\_backtrack\_factor}, & \\
 & \texttt{armijo\_step\_growth} & \\
\bottomrule
\end{tabular}
\end{table}

Everything expensive is shared. The mesh, the land-cover arrays, the carrying capacity, the diffusion tensor and the operator $B$ are built once at start-up and reused by every forward solve, every backward solve and every line-search trial across every outer iteration, and because $B$ is symmetric the two block factorizations of \eqref{eq:blockdiag} serve all of them. The disease-free field $S_\infty$ is shared more widely still: being independent of $v$, it is computed once per mesh bundle and cached on disk (Section~\ref{sec:spinup}), so both the controlled and the uncontrolled driver start from identical data and the two trajectories can be compared directly. Equivalently, the value of $\hat J$ reported at iteration $0$ of a run started from $v\equiv0$ is precisely the do-nothing baseline against which the optimized strategy should be judged.

\section{Discussion and Future Work}

The model presented here extends the classical directly-transmitted wildlife disease framework to chronic wasting disease by explicitly including an environmental prion reservoir $W$, coupled to realistic terrain geometry and land cover heterogeneity. The surface FEM discretization of Section~\ref{sec:fem} exploits the fact that $W$ carries no spatial operator, reducing its update to an independent, exactly-integrable nodal ODE, while $S$, $I$ and $R$ are advanced with an implicit-diffusion, explicit-reaction IMEX scheme. Sections~\ref{sec:control}--\ref{sec:control-implementation} then place a management problem on top of the model, deriving the adjoint system that characterizes an optimal culling strategy and solving it with a forward--backward sweep whose gradient feeds a projected quasi-Newton method.

Future work includes: reconciling the transmission coefficients against the fitted estimates of \citet{Wasserberg2009} and \citet{Miller2006}, which would replace the $R_0$-targeting argument of Section~\ref{sec:params-r0} with a direct calibration; calibrating the land-cover coefficients of Table~\ref{tab:landcover-values} and the carrying capacity $K_0$ against regional surveillance and density data, which are the largest remaining unsupported inputs; calibrating the cost weights against the actual budget and valuation of a management program, $c_4$ in particular having a direct reading as a per-deer removal cost; and exploring spatially varying environmental persistence $\delta(\mathbf{x})$, for example through soil-type dependence. The control formulation itself admits several natural extensions: splitting $R$ into pre-clinical and clinical shedding classes, which is what a sight-based cull would require; additional controls, such as carcass removal (acting as $-v_2 W$ in \eqref{eq:cW}) or non-selective harvest acting on $S$, $I$ and $R$ together; constraints beyond simple bounds, such as a total-budget constraint $\int_0^T\int_\Gamma v \le \mathcal{B}$ or restriction of $v$ to a union of accessible management units; and a terminal payoff penalizing the burden left at $t=T$, which changes only the terminal condition on $\lambda$.

\bibliographystyle{unsrtnat}
\bibliography{references}

\end{document}
